\documentclass[11pt,reqno]{amsart}

\usepackage[T1]{fontenc}
\usepackage[utf8]{inputenc}
\usepackage[margin=1.15in]{geometry}
\usepackage{amsmath,amssymb,amsthm,mathrsfs}
\usepackage{booktabs}
\usepackage{array}
\usepackage{longtable}
\usepackage{enumitem}
\usepackage[hidelinks]{hyperref}

%% ------------------------------------------------------------------
%%  Theorem environments
%% ------------------------------------------------------------------
\theoremstyle{plain}
\newtheorem{theorem}{Theorem}[section]
\newtheorem{proposition}[theorem]{Proposition}
\newtheorem{lemma}[theorem]{Lemma}
\newtheorem{corollary}[theorem]{Corollary}
\newtheorem{conjecture}[theorem]{Conjecture}

\theoremstyle{definition}
\newtheorem{definition}[theorem]{Definition}
\newtheorem{example}[theorem]{Example}
\newtheorem{observation}[theorem]{Observation}

\theoremstyle{remark}
\newtheorem{remark}[theorem]{Remark}
\newtheorem{trap}[theorem]{Trap}

%% ------------------------------------------------------------------
%%  Evidence-class markers.  Every displayed assertion in this paper
%%  carries exactly one of these.
%% ------------------------------------------------------------------
\newcommand{\PROVED}{{\normalfont\sffamily\bfseries[PROVED]}}
\newcommand{\VER}[1]{{\normalfont\sffamily\bfseries[VERIFIED: #1]}}
\newcommand{\EXCL}[1]{{\normalfont\sffamily\bfseries[EXCLUDED: #1]}}
\newcommand{\OPEN}{{\normalfont\sffamily\bfseries[OPEN]}}
\newcommand{\QUOTEDMARK}{{\normalfont\sffamily\bfseries[QUOTED]}}
\newcommand{\QUOTED}[1]{\QUOTEDMARK\,#1}

%% ------------------------------------------------------------------
%%  Notation
%% ------------------------------------------------------------------
\newcommand{\ZZ}{\mathbb{Z}}
\newcommand{\QQ}{\mathbb{Q}}
\newcommand{\RR}{\mathbb{R}}
\newcommand{\CC}{\mathbb{C}}
\newcommand{\NN}{\mathbb{N}}
\newcommand{\dd}{\mathrm{d}}
\newcommand{\eps}{\varepsilon}
\newcommand{\dps}{\mathrm{dps}}
\DeclareMathOperator{\Ord}{ord}
\DeclareMathOperator{\Disc}{disc}
\DeclareMathOperator{\Span}{span}
\DeclareMathOperator{\Image}{Image}
\DeclareMathOperator{\Sol}{Sol}
\DeclareMathOperator{\lcm}{lcm}
\DeclareMathOperator{\den}{den}
\newcommand{\Ihat}{\widehat{I}}
\newcommand{\Phat}{\widehat{P}}
\newcommand{\Ltilde}{\widetilde{L}}

\begin{document}

\title[Frobenius constants and the middle-root obstruction]
      {Frobenius constants, Stokes ratios, and the\\ middle-root obstruction for cellular
       approximations to $\zeta(5)$}

\author{[Author placeholder]}
\address{[Address placeholder]}
\email{[email placeholder]}

\date{\today}

\subjclass[2020]{Primary 11J72; Secondary 11M06, 33C20, 34M35, 14D07}
\keywords{Frobenius constants, gamma conjecture, Ap\'ery-like operators, cellular integrals,
odd zeta values, Stokes constants, irrationality}

\begin{abstract}
We compute the Frobenius constants, in the sense of Bloch--Vlasenko, of the order-nine
Picard--Fuchs operator attached to the Brown--Zudilin family of cellular approximations to
$\zeta(5)$, and we use them to describe an obstruction intrinsic to that family. Three
things are new. First, an instrument: for an Ap\'ery-like operator the Frobenius deformation
series $\kappa(\eps)$ equals the ratio $S(\eps)/S(0)$ of Stokes constants of the dominant
Birkhoff normal solution, the connection point and the growth rate cancelling identically.
This replaces a limit converging like $1/n$ by one converging beyond all orders, and delivers
$434$ digits of $\kappa$ in seconds. Reproducing the published Golyshev--Zagier and
Bloch--Vlasenko tables to $219$ digits, we settle two misprints in the latter, stating both the
printed and the corrected values. Second, the $\kappa$-vector of the Brown--Zudilin operator at
its nearest conifold: the identifications $\kappa_2=-4\zeta(2)$, $\kappa_3=\tfrac{550}{87}\zeta(3)$,
$\kappa_4=\tfrac{365}{29}\zeta(4)$,
$\kappa_5=\tfrac{514}{87}\zeta(5)-\tfrac{2200}{87}\zeta(2)\zeta(3)$ --- read off by
integer-relation detection and verified to $434$ digits, not proved --- make the primitive part
of the first higher Frobenius constant the pure value $\lambda_5=\tfrac{514}{87}\zeta(5)$, with
the $\zeta(2)\zeta(3)$ impurity exactly the decomposable term $\lambda_2\lambda_3$; the passage
from the $\kappa_j$ to the $\lambda_j$ is proved, the closed forms themselves are numerical.
Third, an
identification of the family's purity defect as the Tate class $c=-3/\pi^{2}=12/(2\pi i)^{2}$,
certified to $441$ digits, together with closed forms for the connection constants
($A_QA_{I'}A_I=-\pi^{5/2}/(12\sqrt{37})$, each factor a square root in an $S_3$ cubic field
times a half-integer power of $\pi$, with $\Gamma(r/s)$ excluded for every $s$ searched, namely
$s\in\{3,4,5,6,8,12,24\}$). We formulate a
falsifiable conservation law $r=1+\lfloor m/2\rfloor$ relating the recurrence order to the
multiplicity of the exponent $0$, verified on four families including a prediction at weight
$7$; and we apply the resulting picture to show that the weight-$7$ cellular family admits no
pair-worthiness --- under the one standing hypothesis $s,\Phat\in\Sol(L)$, which we state but do
not prove --- its primitive linear form riding the middle characteristic root by a margin of
$4.635$ nats.
\end{abstract}

\maketitle

\setcounter{tocdepth}{1}
\tableofcontents

%% ==================================================================
\section{Introduction}
%% ==================================================================

\subsection{Frobenius constants}
Let $L$ be a linear differential operator with regular singular point at the origin, and let
$\Phi(s,z)=\sum_{n\ge 0}a_n(s)z^{n+s}$ be its Frobenius deformation, normalised by $a_0=1$.
Golyshev and Zagier \cite{GZ1} discovered, in the course of proving the gamma conjecture for
Fano threefolds of Picard rank one, that the Taylor coefficients of a suitably renormalised
limit of $\Phi$ are zeta values: for the Ap\'ery operator attached to $\zeta(3)$ they found
\begin{equation}\label{eq:apery3}
  \kappa_0=1,\qquad \kappa_1=0,\qquad \kappa_2=-2\zeta(2),\qquad \kappa_3=\tfrac{17}{6}\zeta(3).
\end{equation}
Bloch and Vlasenko \cite{BV} gave the invariant definition. If $c\ne 0,\infty$ is a
\emph{special reflection point} of $L$ --- a regular singular point at which $\sigma_c-1$ has
rank-one image on solutions, and at which all $\sigma_c$-invariant solutions of $L^\vee$ are
analytic --- and $\gamma$ is a path from $0$ to $c$ through regular points, then the
\emph{Frobenius constants} $\kappa_{\rho,n}$ are defined by continuing the Frobenius functions
$\varphi_{\rho,n}=\tfrac{1}{n!}\partial_s^n\Phi|_{s=\rho}$ along $\gamma$ and setting
\begin{equation}\label{eq:BVdef}
  (\sigma_c-1)\,\varphi_{\rho,n}(t) \;=\; \kappa_{\rho,n}\,\delta(t),
\end{equation}
where $\delta$ spans the one-dimensional image of $\sigma_c-1$ \QUOTED{\cite[Def.~22]{BV}}. The
normalisation is $\kappa_0=1$; there is no $2\pi i$ and no $\Gamma$-factor in it. Bloch and
Vlasenko prove that these constants are the expansion coefficients of a \emph{gamma function}
attached to $L^\vee$, and that for $L$ of Picard--Fuchs type they lie in the algebra of periods
with $2\pi i$ inverted \QUOTED{\cite[Thm.~30, Cor.~31]{BV}}. Kerr \cite{Kerr} recasts the
construction in terms of unipotent extensions; Roy and Vlasenko \cite{RV} compute it for
families of elliptic curves, where the higher $\kappa$'s become iterated integrals of modular
forms; Golyshev, Kerr and Sasaki \cite{GKS} relate the adjacent notion of Ap\'ery limits to
extension classes. The relation with Golyshev's \emph{Ap\'ery constants} \cite{Gol} is one of
family resemblance only: those are limits $\lim b_n/a_n$, a different invariant, and we shall
keep the two apart.

A structural feature of \eqref{eq:BVdef} governs everything below. If $m=m(\rho)$ is the
multiplicity of $\rho$ as a root of the indicial polynomial, then $\varphi_{\rho,0},\dots,
\varphi_{\rho,m-1}$ are honest solutions of $L$, while $\varphi_{\rho,k}$ for $k\ge m$ solves an
inhomogeneous equation. Bloch and Vlasenko put it as follows \QUOTED{\cite[p.~21]{BV}}:

\begin{quote}
``\dots the Frobenius constants corresponding to actual solutions of $L$ (with $k<m(\rho)$) are
periods of the limiting Hodge structure at $t=0$. However, there is no reason to expect that
the operators $(D-\rho)^jL$ with $j>0$ are geometric. From this point of view, it is surprising
that the higher Frobenius constants in the above examples are periods.''
\end{quote}

The constants $\kappa_0,\dots,\kappa_{m-1}$ are therefore \emph{geometric}, and $\kappa_m$ is
the first \emph{higher} one. For the Ap\'ery $\zeta(3)$ operator $m=3$ and the first higher
constant is $\kappa_3=\tfrac{17}{6}\zeta(3)$ --- a pure rational multiple of $\zeta(3)$, of
weight exactly $3$.

\subsection{What was not known beyond rank two}
All the operators for which Frobenius constants have been tabulated in the literature are of
order $2$, $3$ or $4$ with a maximally unipotent monodromy point at the origin: the Ap\'ery
operators, Beauville's six stable elliptic families \cite{RV}, the $17$ rank-one Fano families
\cite{GZ1}, and the hypergeometric Calabi--Yau operators, for which $\kappa$ has a closed form
in $\Gamma$-quotients \QUOTED{\cite[Prop.~26]{BV}}. In every tabulated case the multiplicity of
the exponent $0$ is at most $4$, so the first higher constant has weight at most $4$, and the
values $\zeta(5),\zeta(7),\dots$ appear only in the tail, as \emph{higher} constants of a
low-weight operator: for instance $\kappa_5=\tfrac{7}{3}\zeta(5)-\tfrac{17}{3}\zeta(2)\zeta(3)$
for the Ap\'ery $\zeta(3)$ operator, and $\kappa_5=\zeta(5)-3\zeta(2)\zeta(3)$ for the Ap\'ery
$\zeta(2)$ one. In both, $\zeta(5)$ arrives mixed with $\zeta(2)\zeta(3)$.

Two questions were consequently open.

\begin{enumerate}[label=(\alph*),leftmargin=2.2em]
\item \emph{Is there an operator whose \textup{first} higher Frobenius constant is a pure
rational multiple of $\zeta(5)$?} A search of the literature turned up none; the slot appeared
to be empty.
\item \emph{Do the Frobenius constants say anything about irrationality constructions?} The
$\kappa$'s are attached to a differential operator, whereas irrationality proofs are driven by
the arithmetic of a recurrence. The two are the same datum viewed twice; but the dictionary had
not been used in this direction.
\end{enumerate}

The second question is worth stating against the background of what is known. Ap\'ery's proof
that $\zeta(3)$ is irrational \cite{Ap79,Be79}, presented in \cite{Fi04}, remains the only such
result for an odd zeta value; nothing is known about any single $\zeta(2k+1)$ with $k\ge2$.
Infinitely many odd zeta values are irrational \cite{Riv,BR01}, at least
$2^{(1-\epsilon)\log s/\log\log s}$ among $\zeta(3),\dots,\zeta(s)$ \cite{FSZ}; at least one of
$\zeta(5),\zeta(7),\zeta(9),\zeta(11)$ is irrational \cite{Zu02a,Zu04}, at least two of
$\zeta(5),\dots,\zeta(35)$ \cite{LZ}, and one of $\zeta(5),\dots,\zeta(25)$ by elementary means
\cite{Zu18}. Structural reformulations have come from Siegel's lemma \cite{Fi21}, from
transfinite diameter \cite{Br26}, and from the motivic viewpoint \cite{Br16}. What all the
constructive proofs share is a recurrence and a saddle; the results below concern which saddle a
\emph{pure} linear form is allowed to occupy.

The natural test object for both questions is the Brown--Zudilin family. Brown \cite{Br16} interpreted
Ap\'ery-type irrationality proofs as period computations on moduli spaces $\mathcal M_{0,n}$;
Brown and Zudilin \cite{BZ} constructed from cellular integrals on $\mathcal M_{0,8}$ an
effective sequence of rational approximations to $\zeta(5)$ with $|\zeta(5)-p/q|<q^{-0.86}$,
the best effective exponent known. Their integrals decompose as
\begin{equation}\label{eq:BZdeco}
\begin{gathered}
  I \;=\; Q\cdot\bigl(2\zeta(5)+4\zeta(3)\zeta(2)\bigr)-4\Phat\cdot\zeta(2)-2P
  \;=\; 2I'+4\,I''\zeta(2),\\
  I'=Q\zeta(5)-P,\qquad I''=Q\zeta(3)-\Phat ,
\end{gathered}
\end{equation}
and $Q,P,\Phat$ satisfy a third-order Ap\'ery-like recurrence with characteristic polynomial
$4\lambda^3-2368\lambda^2-188\lambda+1$ \QUOTED{\cite[\S2]{BZ}}. The whole difficulty of the
family is visible in \eqref{eq:BZdeco}: the \emph{pure} linear forms $I'$ and $I''$ ride the
middle root $\lambda_2$, while the smallest root $\lambda_1$ --- the one that would give
worthiness exceeding $1$, hence irrationality --- is monopolised by the \emph{impure}
combination $I$, whose leading period $\zeta(5)+2\zeta(2)\zeta(3)$ has depth $2$. We call this
the \emph{middle-root obstruction}. The purpose of this paper is to compute the Frobenius data
of the operator behind \eqref{eq:BZdeco}, to identify the constant that measures the
obstruction, and to test whether the obstruction is a feature of the particular family or of
its weight.

\subsection{Results}
\begin{enumerate}[label=\textbf{(\arabic*)},leftmargin=2.6em]

\item \textbf{The instrument} (\S\ref{sec:instrument}). We prove (Theorem \ref{thm:stokes})
that for an Ap\'ery-like operator with a strictly dominant simple characteristic root
$\lambda$, and with $c=1/\lambda$ the nearest singularity,
\[
  \kappa(\eps)\;=\;\frac{S(\eps)}{S(0)},\qquad
  S(\eps):=\lim_{n\to\infty}\frac{a_n(\rho+\eps)}{A(n+\eps)},
\]
where $A$ is the dominant Birkhoff normal solution. The factors $c^{\eps}$ and $\lambda^{\eps}$
of \cite[Lemma 24]{BV} cancel identically. The resulting algorithm converges beyond all orders
and gives $434$ digits where the direct implementation of \cite[Lemma 24]{BV} gives nine.

\item \textbf{Validation and two errata} (\S\ref{sec:validate}). We reproduce the published
$\kappa$- and $\lambda$-tables of \cite{GZ1,GZ2,BV,RV} at $219$ digits and record two misprints
in \cite{BV}, quoting the printed and the corrected value in each case, and give new values
$\kappa_8,\kappa_9,\lambda_8,\lambda_9$ for the Ap\'ery $\zeta(2)$ operator.

\item \textbf{The $\kappa$-vector of the Brown--Zudilin operator, and a pure $\lambda_5$}
(\S\S\ref{sec:BZop}--\ref{sec:kappavec}). The operator has order $9$, with the exponent $0$
occurring with multiplicity $m=5$ at the origin, and all three finite singularities are
conifolds with exponent set $\{0,\dots,7\}\cup\{3/2\}$. At the nearest conifold,
\VER{$434$ digits}
\[
  \kappa_2=-4\zeta(2),\quad \kappa_3=\tfrac{550}{87}\zeta(3),\quad
  \kappa_4=\tfrac{365}{29}\zeta(4),\quad
  \kappa_5=\tfrac{514}{87}\zeta(5)-\tfrac{2200}{87}\zeta(2)\zeta(3),
\]
whence, \PROVED\ given those identifications, $\lambda_5=\tfrac{514}{87}\zeta(5)$, filling the
slot (a). The impurity of $\kappa_5$ is precisely the decomposable term $\lambda_2\lambda_3$.

\item \textbf{The purity defect is Tate} (\S\ref{sec:defect}). The constant $c$ that rotates the
middle-ray pair $(I',I'')$ onto the minimal ray is
$c=-1/(2\zeta(2))=-3/\pi^2=12/(2\pi i)^2$, proved from \eqref{eq:BZdeco} and certified against
exact ladders to $441$ digits.

\item \textbf{A conservation law} (\S\ref{sec:rpc}). We conjecture that for Ap\'ery-like
families the order $r$ of the recurrence and the multiplicity $m$ of the exponent $0$ satisfy
$r=1+\lfloor m/2\rfloor$, the rays being indexed by the depth of the period they carry, and we
verify it on four families, the fourth (weight $7$, $m=7$, $r=4$) being a prediction made before
the computation of $m$.

\item \textbf{Connection constants are $\Gamma(1/2)$-monomials} (\S\ref{sec:conn}). We identify
all four Brown--Zudilin connection constants:
$A_QA_{I'}A_I=-\pi^{5/2}/(12\sqrt{37})$, with each factor a square root in the $S_3$ cubic field
of the singular locus times a half-integer power of $\pi$, and with $\Gamma(r/s)$ excluded for
every denominator searched, $s\in\{3,4,5,6,8,12,24\}$.

\item \textbf{No pair-worthiness at weight $7$} (\S\ref{sec:weight7}). For the weight-$7$
cellular family the primitive form $I'$ satisfies $I'=I+\zeta(2)|I''|$ termwise, a sum of two
positive quantities; hence --- given Theorem \ref{thm:sector}'s standing hypothesis
$s,\Phat\in\Sol(L)$, which \S\ref{sec:residual} does not prove --- it cannot decay faster than
$I''$, and the fast sector is excluded without any cancellation argument. The measured denominator budget is $\kappa=7$, tight, so
$\gamma_{5,7}\le 0.3379$ and the shortfall is $4.635$ nats, a factor $103$ per step.

\end{enumerate}

\subsection{Evidence classes and conventions}\label{sec:conventions}
Because a substantial part of what follows is high-precision computation, every assertion in
this paper carries exactly one of the following markers.

\begin{itemize}[leftmargin=2.2em]
\item \PROVED\ --- a complete derivation is given, or the statement is an exact symbolic
identity verified in exact arithmetic (integers, rationals, or polynomial algebra).
\item \VER{$d$ digits}\ --- a high-precision computation agreeing with the stated closed form to
$d$ decimal digits. For an integer-relation identification we state in addition the detection
tolerance and the coefficient bound within which the relation was searched.
\item \EXCL{tol, bound} --- an integer-relation search returned no relation at the stated
tolerance and coefficient bound. An exclusion is a statement about a finite search box and
nothing more.
\item \QUOTEDMARK\ --- a statement copied from a source, with the location; we have re-read the
source at that location.
\item \OPEN\ --- not settled here.
\end{itemize}

Numerical agreement is never presented as proof. Where a quantity is quoted to $d$ digits, $d$
is the observed self-agreement between two independent runs (differing in the number of terms
$n$, in the truncation order, or in the working precision), not the working precision itself.

We write $\theta=z\,\dd/\dd z$, and an Ap\'ery-like operator in the form
$L=\sum_{j=0}^{d}z^{j}q_j(\theta)$ with $q_j\in\QQ[x]$; the indicial polynomial at $z=0$ is
$I(s)=q_0(s)$, and the Frobenius recursion reads
$\sum_{j=0}^{n}q_j(n+s-j)\,a_{n-j}(s)=0$ with $a_0=1$. We write $\zeta_2=\zeta(2)$ where space
demands. Following \cite{BZ} we write $I''$ for the $\zeta(3)$-form; the notation $\Ihat$ is
used interchangeably.

\begin{remark}[Methodology; flagged]\label{rem:method}
The computations reported here were carried out with substantial machine assistance, including
automated search over candidate closed forms. Every identification quoted below was re-derived
and re-checked independently of the search that proposed it, and the precision, tolerance and
coefficient bound of each check are stated at the point of use; the exclusions in
\S\ref{sec:exclusions} and \S\ref{sec:conn} record searches that returned nothing, and are
reported with the same care as the positive findings. Appendix \ref{app:record} lists the
computational record.
\end{remark}

%% ==================================================================
\section{The instrument: \texorpdfstring{$\kappa$}{kappa} as a ratio of Stokes constants}
\label{sec:instrument}
%% ==================================================================

\subsection{The computational form of the definition}
We recall the recipe of \cite[Lemma 24]{BV}, in the form in which we shall use it
\QUOTED{\cite[Lemma 24]{BV}}: suppose the special reflection point $t=c$ is the singularity
nearest $0$, that $c\in\RR_{>0}$, that $|\varphi_{\rho,0}(t)|\to\infty$ as $t\to c^-$, that
$a_n(\rho)$ is eventually of one sign, and that
$\lambda_k:=\lim_n a_n^{(k)}(\rho)/(k!\,a_n(\rho))$ exists. Then $\kappa_{\rho,0}\ne 0$ and
\begin{equation}\label{eq:BVlemma24}
  \frac{\kappa_{\rho,k}}{\kappa_{\rho,0}}\;=\;\sum_{j=0}^{k}\frac{(\log c)^j}{j!}\,\lambda_{k-j}.
\end{equation}
Equivalently, with $\Lambda(\eps):=\lim_{n\to\infty}a_n(\rho+\eps)/a_n(\rho)$,
\begin{equation}\label{eq:kappaLambda}
  \kappa(\eps)\;=\;\sum_{k\ge 0}\kappa_{\rho,k}\eps^k\;=\;c^{\eps}\,\Lambda(\eps).
\end{equation}
This is also, verbatim, the recipe of \cite[\S9]{GZ1}, there written
$\kappa_0(\eps)=C^{-\eps}\lim_n A_n(\eps)/A_n$ with $C=1/c$ the exponential growth rate.

Taken literally, \eqref{eq:kappaLambda} is a poor algorithm: $a_n(\rho+\eps)/a_n(\rho)$
converges like $1/n$, and Richardson extrapolation on it reaches about nine digits before the
conditioning of the extrapolation defeats the gain. Golyshev and Zagier already observed
\QUOTED{\cite[\S9]{GZ1}} that dividing instead by the Birkhoff--Trjitzinsky asymptotic
$A(n+\eps)$, rather than by $a_n(\rho)$, converges faster than any power of $n$, and reported
$300$ digits from $n=100$ in the Ap\'ery case. Theorem \ref{thm:stokes} explains why, and shows
that the correction is not a numerical trick but an identity: the two exponential factors in
\eqref{eq:kappaLambda} cancel exactly, leaving $\kappa$ as a pure ratio of Stokes constants.

\subsection{The Birkhoff normal solution}\label{sec:birkhoff}
Fix an operator $L=\sum_{j}z^{j}q_j(\theta)$ of order $r$ and degree $d$ in $z$, and put
$D=\max_j\deg q_j$. The characteristic polynomial of the associated recurrence is
\[
  \chi(\lambda)=\lambda^{d}\sum_{j}\lambda^{-j}\,[x^{D}]q_j(x),
\]
whose roots are the reciprocals of the finite singularities of $L$. Assume one root $\lambda$ is
simple and strictly dominant in absolute value.

\begin{proposition}[Birkhoff normal solution]\label{prop:birkhoff} \PROVED\
There is a unique formal expression
\begin{equation}\label{eq:Aform}
  A(x)\;=\;\lambda^{x}\,x^{\alpha}\,F(1/x),\qquad F(u)=\sum_{k\ge 0}c_ku^{k},\quad c_0=1,
\end{equation}
satisfying $\sum_{j}q_j(x-j)A(x-j)=0$ as a formal series in $u=1/x$. Writing
$Q_j(u):=u^{D}q_j(1/u-j)\in\QQ[u]$ and
\[
  \beta:=\sum_j\lambda^{-j}Q_j'(0),\qquad
  \gamma:=\sum_j\lambda^{-j}\,j\,Q_j(0),\qquad
  H_k(u):=\sum_j\lambda^{-j}Q_j(u)(1-ju)^{\alpha-k},
\]
one has $\sum_j\lambda^{-j}Q_j(0)=0$ (the characteristic equation),
\begin{equation}\label{eq:alpha}
  \alpha=\beta/\gamma,\qquad\text{and}\qquad
  c_m=-\frac{1}{\gamma m}\sum_{k<m}c_k\,h_{k,\,m+1-k},\qquad h_{k,i}:=[u^{i}]H_k(u).
\end{equation}
\end{proposition}

\begin{proof}
Substituting \eqref{eq:Aform} into $\sum_jq_j(x-j)A(x-j)=0$, dividing by $\lambda^{x}x^{\alpha}$
and multiplying by $u^{D}$ with $u=1/x$ turns the equation into
$\sum_{j}\lambda^{-j}Q_j(u)(1-ju)^{\alpha}F\bigl(u/(1-ju)\bigr)=0$, since
$q_j(x-j)=u^{-D}Q_j(u)$ and $1/(x-j)=u/(1-ju)$. Expanding
$F(u/(1-ju))=\sum_k c_ku^{k}(1-ju)^{-k}$ gives
$\sum_{k\ge0}c_k\,u^{k}H_k(u)=0$. The coefficient of $u^{0}$ is $c_0H_0(0)$, and
$H_k(0)=\sum_j\lambda^{-j}Q_j(0)$ is independent of $k$; so this coefficient vanishes precisely
when $\lambda$ solves the characteristic equation, and then $h_{k,0}=0$ for all $k$. The
coefficient of $u^{1}$ is therefore $c_0h_{0,1}=\beta-\alpha\gamma$, which forces
$\alpha=\beta/\gamma$. For $m\ge 1$ the coefficient of $u^{m+1}$ is
$\sum_{k\le m}c_kh_{k,m+1-k}$, whose top term is
$c_mh_{m,1}=c_m(\beta-(\alpha-m)\gamma)=c_m\,m\gamma$; solving for $c_m$ gives \eqref{eq:alpha}.
Simplicity of $\lambda$ gives $\gamma\ne0$, so the recursion determines all $c_m$ uniquely.
\end{proof}

\begin{remark}\label{rem:alphacheck}
The exponent $\alpha$ comes out an \emph{exact rational} and is a strong self-check on the
implementation: $\alpha=-1$ for the Ap\'ery $\zeta(2)$ operator, $\alpha=-3/2$ for the Ap\'ery
$\zeta(3)$ operator, $\alpha=-5/2$ for the Brown--Zudilin operator. In each case
$\alpha=-\rho-1$ with $\rho$ the non-integer local exponent at the conifold, i.e.\ $\rho=0$,
$1/2$, $3/2$ respectively. \PROVED\ (exact rational arithmetic; for the Brown--Zudilin case the
value $\rho=3/2$ is proved independently in Proposition \ref{prop:conifold}).
\end{remark}

\subsection{The Stokes-ratio identity}

\begin{theorem}[Stokes-ratio form of $\kappa$]\label{thm:stokes} \PROVED\
Let $L$ be as above, with $\lambda$ a simple strictly dominant characteristic root and $A$ the
Birkhoff normal solution of Proposition \ref{prop:birkhoff}. Let $\rho$ be a local exponent at
$z=0$ for which the limit
\[
  S(\eps)\;:=\;\lim_{n\to\infty}\frac{a_n(\rho+\eps)}{A(n+\eps)}
\]
exists and is non-zero for $\eps$ in a neighbourhood of $0$. Let $c=1/\lambda$ be the nearest
singularity, assumed a special reflection point, and suppose $\kappa(\eps)=c^{\eps}\Lambda(\eps)$
as in \eqref{eq:kappaLambda}. Then
\begin{equation}\label{eq:stokes}
  \boxed{\;\kappa(\eps)\;=\;\frac{S(\eps)}{S(0)}\;}
\end{equation}
identically as formal power series in $\eps$.
\end{theorem}

\begin{proof}
The Frobenius deformation acts on the recursion by the substitution $n\mapsto n+s$: the sequence
$\bigl(a_n(\rho+\eps)\bigr)_n$ satisfies $\sum_j q_j\bigl((n+\eps)-j\bigr)a_{n-j}(\rho+\eps)=0$
when $\rho=0$, and in general the same equation with $n+\rho+\eps$ in place of $n+\eps$. Hence
the \emph{same} formal solution $A$ of Proposition \ref{prop:birkhoff}, evaluated at the shifted
argument $x=n+\rho+\eps$, governs its asymptotics; absorbing $\rho$ into the normalisation we
write $A(n+\eps)$. Therefore
\[
  \Lambda(\eps)=\lim_{n\to\infty}\frac{a_n(\rho+\eps)}{a_n(\rho)}
  =\frac{S(\eps)}{S(0)}\cdot\lim_{n\to\infty}\frac{A(n+\eps)}{A(n)} .
\]
By \eqref{eq:Aform},
$A(n+\eps)/A(n)=\lambda^{\eps}\,(1+\eps/n)^{\alpha}\,F(1/(n+\eps))/F(1/n)\to\lambda^{\eps}$.
Thus $\Lambda(\eps)=\lambda^{\eps}S(\eps)/S(0)$ and
$\kappa(\eps)=c^{\eps}\Lambda(\eps)=\lambda^{-\eps}\lambda^{\eps}S(\eps)/S(0)=S(\eps)/S(0)$.
\end{proof}

\begin{remark}[Why this is fast]\label{rem:fast}
The error in $a_n(\eps)/A(n+\eps)-S(\eps)$ has two parts: the contribution of the subdominant
solutions, of size $O\bigl((\lambda_{\mathrm{next}}/\lambda)^{n}\bigr)$, and the truncation of
the divergent series $F$ after $M$ terms, of size $\approx M!/(\mathcal S n)^{M}$ with
$\mathcal S=\log|\lambda/\lambda_{\mathrm{next}}|$. Both are governed by the same $\mathcal S$,
and optimal truncation $M\approx \mathcal S n$ gives error $e^{-\mathcal S n}$: the convergence
is beyond all orders. For the Ap\'ery $\zeta(3)$ operator
$\mathcal S=2\log(17+12\sqrt2)=7.05099\ldots$, so $n=100$ yields
$7.05099\cdot 100/\log 10=306.2$ digits --- reproducing exactly the ``$300$ digits from
$n=100$'' of \cite[\S9]{GZ1}. \PROVED\ (as an asymptotic count; the observed precisions in
\S\ref{sec:validate} and \S\ref{sec:kappavec} agree with the model, see Table \ref{tab:precision}.)
\end{remark}

\begin{table}[t]
\centering
\caption{The truncation model of Remark \ref{rem:fast} against observed self-agreement.
$\mathcal S=\log|\lambda/\lambda_{\mathrm{next}}|$; predicted $\log_{10}$ error is
$\log_{10}\bigl(M!/(\mathcal S n)^{M}\bigr)$. \VER{$208$--$434$ digits, as tabulated}}
\label{tab:precision}
\begin{tabular}{lccccr}
\toprule
operator & $\mathcal S$ & $n$ & $M$ & predicted digits & observed \\
\midrule
Ap\'ery $\zeta(3)$ & $7.0510$ & $300$ & $130$ & $212$ & $208$--$210$ \\
Ap\'ery $\zeta(3)$ & $7.0510$ & $400$ & $130$ & $229$ & (capped by $\dps=220$) \\
Brown--Zudilin     & $8.8560$ & $700$ & $260$ & $469$ & $434$ \\
Brown--Zudilin     & $8.8560$ & $800$ & $260$ & $485$ & (capped by $\dps=520$) \\
\bottomrule
\end{tabular}
\end{table}

\subsection{Relation to the $p$-adic side, stated without overreach}\label{sec:padic}
There is a $p$-adic story running parallel to this one, and it is easy to overstate the
parallel. On the $p$-adic side, the limits of Frobenius towers attached to the same families
are \emph{values} of $G(x)=\Gamma_p(x)e^{-\Gamma_p'(0)x}=\exp\bigl(-\sum_{m\ge2}\zeta_p(m)x^m/m\bigr)$
at $p$-power arguments, whereas the Frobenius \emph{structure} constants of
\cite{BeVl25} are the Taylor coefficients of the same $G$ at $0$. What is
verified here, and all that we assert, is:
\begin{enumerate}[label=(\alph*),leftmargin=2.2em]
\item the archimedean series $\log\kappa(\eps)=\sum_j\lambda_j\eps^{j}$ has
$\lambda_j\in\QQ\cdot\zeta(j)$ for $j\le5$ on the Brown--Zudilin and Ap\'ery $\zeta(3)$
operators, and for $j\le4$ on the Ap\'ery $\zeta(2)$ operator --- an initial segment of exactly
the shape $\exp\bigl(-\sum_m r_m\zeta(m)x^m\bigr)$, $r_m\in\QQ$;
\item the archimedean connection constants are $\Gamma(1/2)$-monomials times algebraic numbers
(\S\ref{sec:conn}).
\end{enumerate}
We did \emph{not} find, and do not claim, a literal identity between the two $\Gamma$'s. The
honest statement is that both sides are expansion data of a $\Gamma$-type generating function
attached to the operator --- $\Gamma_p$ on one side, the $\Gamma_{\xi_0}$ of
\QUOTED{\cite[Thm.~30]{BV}} on the other --- and that in both cases the coefficients are
(multiple) zeta values.

%% ==================================================================
\section{Validation, and two errata}\label{sec:validate}
%% ==================================================================

Reproducing the published values is the credibility of everything that follows, and we report it
in full.

\subsection{Controls}
Run parameters: working precision $\dps=220$, $\eps$-truncation $K=13$, Birkhoff truncation
$M=130$, $n\in\{300,400\}$; self-agreement across $n$ is $208$--$210$ digits, in agreement with
the model of Remark \ref{rem:fast} (Table \ref{tab:precision}).

\begin{table}[t]
\centering
\caption{Instrument validation. All residuals are the difference between the computed value and
the published closed form; each source location was re-read \QUOTEDMARK. \VER{219 digits}}
\label{tab:validate}
\begin{tabular}{@{}p{0.30\textwidth}p{0.16\textwidth}p{0.46\textwidth}@{}}
\toprule
control & source & result \\
\midrule
Ap\'ery $\zeta(3)$: exponent $\alpha$ & --- &
  $-3/2$, exact rational \\[3pt]
Ap\'ery $\zeta(3)$: Stokes constant $S(0)$ & classical &
  $0.2200437671126430378506897598\ldots$, equal to
  $(1+\sqrt2)^2\bigl/\bigl(2^{9/4}\pi^{3/2}\bigr)$ to all digits computed \\[3pt]
%% FIXED [MINOR]: "the two exceptions" sat on the Apery-zeta(2) row, but Erratum
%% FIXED \ref{obs:err1} belongs to the Apery-zeta(3) row (BV Ex.~29) and Erratum \ref{obs:err2}
%% FIXED to the zeta(2) row (BV Ex.~28). One each.
Ap\'ery $\zeta(3)$: $\kappa_2,\kappa_3,\kappa_4$ & \cite[Ex.~29]{BV} &
  residuals $<10^{-219}$ against $-2\zeta(2)$, $\tfrac{17}{6}\zeta(3)$, $2\zeta(4)$; for
  $\kappa_5$ see Erratum \ref{obs:err1} \\[3pt]
Ap\'ery $\zeta(3)$: $\kappa_6,\dots,\kappa_{10}$ & \cite[(47)]{GZ2} &
  residuals $<10^{-219}$ \\[3pt]
Ap\'ery $\zeta(3)$: $\lambda_2,\dots,\lambda_{10}$ & \cite[\S9]{GZ1} &
  all reproduced; the two published packagings are mutually consistent, see Remark
  \ref{rem:packaging} \\[3pt]
Ap\'ery $\zeta(2)$: $\kappa_2,\dots,\kappa_7$ &
  \cite[Ex.~28]{BV}, \cite[Tab.~2 D]{RV} & residuals $<10^{-219}$, with the one
  exception recorded in Erratum \ref{obs:err2} \\
\bottomrule
\end{tabular}
\end{table}

\begin{remark}[The two published packagings agree]\label{rem:packaging} \VER{60 digits}
The tables of \cite{GZ1} and \cite{GZ2} present the same data twice: \cite[(47)]{GZ2} lists
$\kappa_j$ for $j\le11$, \cite[\S9]{GZ1} lists $\lambda_j=[\eps^j]\log\kappa(\eps)$ for the same
range. Exponentiating the $\lambda$-list reproduces the $\kappa$-list, and vice versa, to $60$
digits --- the limit here being the number of digits carried by the printed rational
coefficients, not the computation. Two re-packagings we record because they are convenient and
are not printed in that form: $\lambda_6=-\tfrac{2}{3}\zeta(6)-\tfrac{1}{72}\zeta(3)^2
=-\tfrac{16}{105}\zeta(2)^3-\tfrac{1}{72}\zeta(3)^2$ and
$\lambda_8=\tfrac{29}{12}\zeta(8)-\tfrac{11}{18}\zeta(3)\zeta(5)
=\tfrac{58}{175}\zeta(2)^4-\tfrac{11}{18}\zeta(3)\zeta(5)$. \PROVED\ (exact rational identities
$\zeta(6)=\tfrac{8}{35}\zeta(2)^3$, $\zeta(8)=\tfrac{24}{175}\zeta(2)^4$.)
\end{remark}

\subsection{Two errata in \texorpdfstring{\cite{BV}}{[BV]}}\label{sec:errata}
Both printed forms below were read verbatim from the published text before the computation was
run; both corrections agree with \cite[(47)]{GZ2} and \cite[Table 2]{RV} respectively. We state
them plainly and with no implication beyond a typesetting slip.

\begin{observation}[Erratum 1]\label{obs:err1} \VER{219 digits}
\cite[Example 29]{BV} prints, for the Ap\'ery $\zeta(3)$ operator,
\[
  \kappa_5=\tfrac{7}{5}\zeta(5)-\tfrac{17}{3}\zeta(2)\zeta(3)
  \qquad\text{\QUOTED{\cite[Ex.~29]{BV}}}.
\]
The computed value is
\[
  \kappa_5=-8.78522655635014817501029803334929851633233616556963\ldots
\]
The form $\tfrac{7}{3}\zeta(5)-\tfrac{17}{3}\zeta(2)\zeta(3)$ agrees with it to
$4.6\times10^{-219}$; the printed form differs by $0.9678$. The correct value is therefore
\[
  \kappa_5=\tfrac{7}{3}\zeta(5)-\tfrac{17}{3}\zeta(2)\zeta(3),
\]
in agreement with $\kappa_5=-\tfrac{17}{18}\pi^2\zeta(3)+\tfrac{7}{3}\zeta(5)$ of
\QUOTED{\cite[(47)]{GZ2}}. The erratum concerns the coefficient of $\zeta(5)$ only; the term
$-\tfrac{17}{3}\zeta(2)\zeta(3)$ stands as printed. All other values listed in
\cite[Ex.~29]{BV} check out.
\end{observation}

\begin{observation}[Erratum 2]\label{obs:err2} \VER{219 digits}
\cite[Example 28]{BV} prints, for the Ap\'ery $\zeta(2)$ operator,
\[
  \kappa_6=\tfrac{87}{16}\zeta(6)-\tfrac{5}{2}\zeta(3)^2
  \qquad\text{\QUOTED{\cite[Ex.~28]{BV}}}.
\]
The computed value is
\[
  \kappa_6=9.14415489562452778198190394025559993842547\ldots
\]
The form with $+\tfrac{5}{2}\zeta(3)^2$ agrees to $5.3\times10^{-220}$; the printed form differs
by $7.2247$. The correct value is
\[
  \kappa_6=\tfrac{87}{16}\zeta(6)+\tfrac{5}{2}\zeta(3)^2,
\]
in agreement with \QUOTED{\cite[Table 2, case D]{RV}}. The remaining entries
$\kappa_2,\kappa_3,\kappa_4,\kappa_5,\kappa_7$ of \cite[Ex.~28]{BV} check out.
\end{observation}

\subsection{New values for the Ap\'ery \texorpdfstring{$\zeta(2)$}{zeta(2)} operator}
The published tables for this operator stop at $\kappa_7$. We add
\begin{align}
  \kappa_8 &= \tfrac{8627}{4200}\zeta(2)^4-\tfrac{7}{2}\zeta(2)\zeta(3)^2
              +\tfrac{15}{2}\zeta(3)\zeta(5),\notag\\
  \kappa_9 &= \tfrac{29}{21}\zeta(2)^3\zeta(3)+\tfrac{5}{3}\zeta(3)^3
              -\tfrac{19}{4}\zeta(2)^2\zeta(5)+\tfrac{115}{16}\zeta(2)\zeta(7)
              -\tfrac{2011}{72}\zeta(9),\label{eq:apery2new}\\
  \lambda_8 &= \tfrac{336593}{105000}\zeta(2)^4+\tfrac{2}{5}\zeta(2)\zeta(3)^2
              +\tfrac{11}{2}\zeta(3)\zeta(5),\notag\\
  \lambda_9 &= -\tfrac{14263}{5250}\zeta(2)^3\zeta(3)-\tfrac{2}{3}\zeta(3)^3
              -\tfrac{649}{100}\zeta(2)^2\zeta(5)-\tfrac{39}{16}\zeta(2)\zeta(7)
              -\tfrac{2011}{72}\zeta(9).\notag
\end{align}
\VER{219 digits}. As an internal check, the $\lambda$-values were re-derived from the
$\kappa$-values by exponentiation --- an exact algebraic consequence of the
$\kappa$-identifications \PROVED\ --- and agree numerically to the full precision of the printed
rational data ($60$ digits). In the same
spirit, $\lambda_4=-\tfrac{39}{20}\zeta(4)$ and
$\lambda_5=\zeta(5)-\tfrac{1}{5}\zeta(2)\zeta(3)$ for this operator, both exact consequences of
$\kappa_2=-\tfrac75\zeta(2)$, $\kappa_3=2\zeta(3)$, $\kappa_4=\tfrac12\zeta(4)$,
$\kappa_5=\zeta(5)-3\zeta(2)\zeta(3)$ \PROVED. We did not attempt $\kappa_{10}$, which requires
$\zeta(3,7)$ and $\zeta(2)\zeta(3,5)$ and hence a validated high-precision multiple-zeta-value
routine \OPEN.

%% ==================================================================
\section{The Brown--Zudilin operator}\label{sec:BZop}
%% ==================================================================

\subsection{From the recurrence to the operator}
Brown and Zudilin give the order-three recurrence satisfied by $I_n$, hence by $Q_n$, $P_n$,
$\Phat_n$, obtained by creative telescoping from the cellular integrals evaluated with
%% FIXED [MINOR]: bibitem \cite{Zu02} was uncited; it is Zudilin's third-order Apery-like
%% FIXED recursion for zeta(5), i.e. the precursor of \eqref{eq:BZrec}, and is cited here.
\cite{Pa15} and \cite{Ko10} \QUOTED{\cite[\S2]{BZ}}; a third-order Ap\'ery-like recursion for
$\zeta(5)$ of this shape goes back to \cite{Zu02}:
\begin{equation}\label{eq:BZrec}
  A(n)u_{n+1}+B(n)u_n+C(n)u_{n-1}+D(n)u_{n-2}=0,
\end{equation}
with
\begin{align*}
  A(n)&=2(2n+1)(41218n^3-48459n^2+20010n-2871)(n+1)^5,\\
  B(n)&=-\bigl(97604224n^9+178061760n^8+72005308n^7-48634688n^6-39076836n^5\\
      &\qquad +2622730n^4+7581006n^3+920112n^2-543402n-120582\bigr),\\
  C(n)&=-2n\bigl(3874492n^8-2617900n^7-3144314n^6+2947148n^5+647130n^4\\
      &\qquad -1182926n^3+115771n^2+170716n-44541\bigr),\\
  D(n)&=n(41218n^3+75195n^2+46746n+9898)(n-1)^5 .
\end{align*}
With $\theta=z\,\dd/\dd z$ and $Y(z)=\sum_nu_nz^n$ this transports to
\begin{equation}\label{eq:BZop}
  L\;=\;A(\theta-1)+z\,B(\theta)+z^2\,C(\theta+1)+z^3\,D(\theta+2),
\end{equation}
of order $9$ in $\theta$ and degree $3$ in $z$. \PROVED\ (exact symbolic computation). The
residual of \eqref{eq:BZrec} is exactly $0$ for $u\in\{Q,P,\Phat\}$ at every $n=2,\dots,359$
($1074$ exact rational checks), with anchors matching the printed values
$Q=1,21,2989$; $P=0,\tfrac{87}{4},\tfrac{1190161}{384}$;
$\Phat=0,\tfrac{101}{4},\tfrac{344923}{96}$ of \QUOTED{\cite[\S2]{BZ}}.
\VER{exact rational arithmetic, $1074$ checks; no digits are at issue}. It also
holds at $n=1$ for $Q$ and $\Phat$ but not for $P$, which carries an inhomogeneity there.

\subsection{Local data}

\begin{proposition}\label{prop:local} \PROVED\ (exact symbolic computation.)
For the operator \eqref{eq:BZop}:
\begin{enumerate}[label=(\roman*),leftmargin=2.2em]
\item the coefficient of $\theta^{9}$ is $41218\,(z^3-188z^2-2368z+4)$, the reciprocal of the
characteristic polynomial $4\lambda^3-2368\lambda^2-188\lambda+1$; the finite nonzero
singularities are $z_i=1/\lambda_i$;
\item the indicial polynomial at $z=0$ is
\[
  I(s)=q_0(s)=2\,s^{5}(2s-1)\bigl(41218s^3-172113s^2+240582s-112558\bigr),
\]
so the exponent $0$ has multiplicity $m=5$, and the other exponents are $1/2$ and the three
roots of an irreducible cubic;
\item the $z^3$ part is $(\theta+1)^5(\theta+2)(41218\theta^3+322503\theta^2+842142\theta+733914)$,
the same fivefold structure at $z=\infty$;
\item $\Disc(z^3-188z^2-2368z+4)=2^4\cdot 37^3\cdot 557^2$, not a square, so the cubic field
$\QQ(\lambda_3)$ has Galois group $S_3$; note $41218=2\cdot37\cdot557$.
\end{enumerate}
\end{proposition}

The point $z=0$ is thus \emph{not} a MUM point for the order-$9$ operator --- that would require
multiplicity $9$ --- but it carries a rank-$5$ maximally unipotent sub-block, with local
logarithmic depth up to $\log^4z$. The multiplicity equals the motivic weight,
$m=5$. This resolves an apparent tension in the literature-adjacent folklore: the
\emph{recurrence} \eqref{eq:BZrec} has order $3$, and one might infer a local log-depth of at
most $2$, too little room for weight $5$; but the \emph{operator} has a rank-$5$ unipotent block,
which is exactly enough.

\begin{proposition}[All three finite singularities are conifolds]\label{prop:conifold} \PROVED\
Let $c_9,c_8$ be the coefficients of $\theta^9,\theta^8$ in \eqref{eq:BZop}. At a simple root
$z_0$ of $c_9$ the local exponents of an order-$9$ operator are
$\{0,1,\dots,7\}\cup\{\rho\}$ with $\rho=8-c_8(z_0)/\bigl(z_0c_9'(z_0)\bigr)$. Here
\[
  \rho(z)=\frac{-679z^3+106784z^2+1082176z-1384}{74z(-3z^2+376z+2368)},
\]
and $\rho(z)-\tfrac32$ has numerator $-173\,(z^3-188z^2-2368z+4)$, a nonzero rational multiple
of the singular cubic. Hence $\rho\equiv\tfrac32$ identically on the singular locus:
\[
  \text{all three singularities }z_1,z_2,z_3\text{ are conifold points with exponents }
  \{0,1,\dots,7\}\cup\{3/2\}.
\]
\end{proposition}

This is an exact algebraic identity, not a numerical fit, and it has a structural payoff:
exactly one non-integer exponent at each singular point means that $\sigma_c-1$ has rank-one
image there, so the first condition of \QUOTED{\cite[Def.~21]{BV}} holds at all three
singularities. It also confirms $\alpha=-\rho-1=-5/2$ for every ray, matching the measured
value of Remark \ref{rem:alphacheck}.

\subsection{Applicability of \texorpdfstring{\cite[Def.~22]{BV}}{[BV, Def. 22]}, and one
hypothesis that fails}\label{sec:applic}
We record precisely what holds and what does not.

\begin{itemize}[leftmargin=1.6em]
\item \emph{Reflection point, first condition:} this holds at all three of the singularities,
by Proposition \ref{prop:conifold}. \PROVED
\item \emph{Reflection point, second condition,} namely analyticity at $c$ of the
$\sigma_c$-invariant solutions of $L^\vee$: not verified here. \OPEN
\item \emph{Nearest singularity real and positive:} one has
$z_3=1/\lambda_3=0.0016889627\ldots$, so the frame of \eqref{eq:BVlemma24} applies. \PROVED
\item \emph{$a_n(0)$ of one sign:} $a_n(0)=Q_n=1,21,2989,714549,\dots$, verified exactly from the
Frobenius recursion --- which is simultaneously the check that the $q_j$ are correct.
\VER{exact, $n\le359$}
\item \emph{Hypothesis $|\varphi_{\rho,0}(t)|\to\infty$ as $t\to c^-$: fails.} With
$\alpha=-5/2$ one has $\sum_nQ_nz_3^{\,n}\asymp\sum n^{-5/2}<\infty$. It \emph{also} fails for
the Ap\'ery $\zeta(3)$ operator ($\alpha=-3/2$), for which \cite{BV} nonetheless quote the
conclusion and for which we verify it to $219$ digits in \S\ref{sec:validate}.
\end{itemize}

\begin{observation}\label{obs:sufficient} \VER{219 digits at $\alpha=-3/2$; consistency checks
at $\alpha=-5/2$, see \S\ref{sec:kappavec}}
Empirically, hypothesis $|\varphi_{\rho,0}(t)|\to\infty$ of \cite[Lemma~24]{BV} is
\emph{sufficient but not necessary}: the conclusion $\kappa(\eps)=c^{\eps}\Lambda(\eps)$ holds at
$\alpha=-3/2$ and at $\alpha=-5/2$. This is the licence under which we run the machinery on the
Brown--Zudilin operator. It is an empirical licence, and we flag it as such.
\end{observation}

Finally, the rank-$5$ block is a feature and not an obstruction. Definition 22 of \cite{BV}
never requires MUM at $z=0$; it requires a reflection point. What $m(0)=5$ does is decide
\emph{which} $\kappa_j$ is the first higher one: $\kappa_0,\dots,\kappa_4$ are periods of the
limiting mixed Hodge structure, and $\kappa_5$ is the first non-geometric constant --- of weight
$5$. That is exactly why this operator is the right place to look for the $\zeta(5)$-analogue of
$\kappa_3=\tfrac{17}{6}\zeta(3)$.

%% ==================================================================
\section{The \texorpdfstring{$\kappa$}{kappa}-vector, and a pure primitive
\texorpdfstring{$\lambda_5$}{lambda 5}}\label{sec:kappavec}
%% ==================================================================

\subsection{The computation}
Run parameters: $\dps=520$, $\eps$-truncation $K=16$, Birkhoff truncation $M=260$,
$n\in\{700,800\}$; self-agreement across $n$ is $434$ digits, consistent with the truncation
model (Table \ref{tab:precision}) and capped by the working precision. As a cross-check of the
entire pipeline, the Stokes constant
\[
  S(0)=0.066676425727156767841653340639345444469628745592456\ldots
\]
reproduces the connection constant $A_Q$ of \S\ref{sec:conn}, computed by Neville extrapolation
in an entirely separate implementation, in every digit of the latter's $158$-digit stability
window. \VER{158 digits, independent implementations}

\begin{theorem}[The Brown--Zudilin $\kappa$-vector]\label{thm:kappavec}
%% FIXED [MINOR x2]: (a) the quoted residuals 1e-440..1e-451 are NOT reproducible from the
%% FIXED archived data -- work/gamma/t3e_high.py:88 writes kappas_hi.json at nstr(x,420), so any
%% FIXED recheck floors at ~1e-419 (referee's and my own reruns both land at ~2e-419). The tag now
%% FIXED says where the deeper residuals live. (b) the printed digit strings are truncated to 48
%% FIXED decimals WITH the last digit rounded (kappa_4,5,7,8 round up), so a bare \ldots was
%% FIXED misleading; stated. (c) \small was being issued inside math mode; moved outside.
\VER{434 digits; integer-relation detection at tolerance $10^{-380}$ with coefficient bound
$10^{14}$; residuals $10^{-440}$ to $10^{-451}$ in the working session, which the archived
$420$-digit dump cannot reproduce below $\approx10^{-419}$}
At the nearest conifold $z_3=1/\lambda_3$, with $\kappa_0=1$ and $\kappa_1=0$
($|\kappa_1|<10^{-440}$), the computed values are, to $48$ decimals with the final digit
rounded,
{\small
\[
\begin{aligned}
  \kappa_2 &= \phantom{-00}{-6.579736267392905745889660666584100756875799604827}\ldots,\\
  \kappa_3 &= \phantom{-000}{7.599210307330768470917884929095373504836120238934}\ldots,\\
  \kappa_4 &= \phantom{-00}{13.622344148433291031149701697845733948716693015529}\ldots,\\
  \kappa_5 &= \phantom{-0}{-43.874582810463900205738690933659318347119467057797}\ldots,\\
  \kappa_6 &= \phantom{-00}{34.702212286332465372786749212708512062179423906148}\ldots,\\
  \kappa_7 &= \phantom{-00}{51.859399033579894576098101769427996696229062468954}\ldots,\\
  \kappa_8 &= {-193.486228134049501553393657261757187130330010383354}\ldots,\\
  \kappa_9 &= \phantom{-}{328.247447939140704375272593880584060942588691516576}\ldots,
\end{aligned}
\]}
and the first four are identified as
\[
  \kappa_2=-4\zeta(2),\qquad
  \kappa_3=\tfrac{550}{87}\zeta(3),\qquad
  \kappa_4=\tfrac{146}{29}\zeta(2)^2=\tfrac{365}{29}\zeta(4),\qquad
  \kappa_5=\tfrac{514}{87}\zeta(5)-\tfrac{2200}{87}\zeta(2)\zeta(3).
\]
\end{theorem}

The weight-$j$ basis used for the identification is
$\{\zeta(2)^a\prod\zeta(\text{odd})\}$, checked for internal relations \emph{before} any
relation was read off; see Trap \ref{trap:basis}.

\subsection{The primitive normalisation}
Write $\lambda(\eps)=\log\kappa(\eps)=\sum_{j\ge2}\lambda_j\eps^{j}$; we call the $\lambda_j$ the
\emph{primitive} Frobenius constants. This is the packaging used in \cite[\S9]{GZ1}.

%% FIXED [CRITICAL]: the \PROVED tag was unsound. Only the RELATION lambda_5 = kappa_5 -
%% FIXED kappa_2 kappa_3 is proved; the closed forms are inherited from Theorem
%% FIXED \ref{thm:kappavec}, which is \VER{434 digits} (PSLQ) and itself runs under the empirical
%% FIXED licence of Observation \ref{obs:sufficient}. The licensing hypothesis is now IN the
%% FIXED statement ("Assume the identifications of Theorem \ref{thm:kappavec}"), and the marker
%% FIXED is split between the two halves. Same repair applied at Observations
%% FIXED \ref{obs:purityatm}, \ref{obs:cheapplane}, \ref{obs:primitive}, in the abstract and in
%% FIXED Results item (3). The identity itself was re-verified independently from
%% FIXED work/gamma/kappas_hi.json at mp.dps=460: kappa_5 - kappa_2 kappa_3 - (514/87)zeta(5)
%% FIXED = 2.34e-419, i.e. zero to the archive's precision floor. The mathematics is right.
\begin{corollary}[A pure $\zeta(5)$]\label{cor:pure}
\PROVED\ (the relations); \VER{$434$ digits} (the closed forms).
\emph{Assume the identifications of Theorem \textup{\ref{thm:kappavec}}.} Then,
for the Brown--Zudilin operator,
\[
  \lambda_2=-4\zeta(2),\qquad \lambda_3=\tfrac{550}{87}\zeta(3),\qquad
  \lambda_4=-\tfrac{86}{29}\zeta(2)^2=-\tfrac{215}{29}\zeta(4),\qquad
  \boxed{\;\lambda_5=\tfrac{514}{87}\zeta(5)\;}
\]
a pure rational multiple of $\zeta(5)$. Moreover $\kappa_5=\lambda_5+\lambda_2\lambda_3$
identically, i.e.\ the entire $\zeta(2)\zeta(3)$ impurity of $\kappa_5$ is the
\emph{decomposable} term
$\lambda_2\lambda_3=(-4\zeta(2))\bigl(\tfrac{550}{87}\zeta(3)\bigr)
=-\tfrac{2200}{87}\zeta(2)\zeta(3)$.
\end{corollary}

\begin{proof}
$\lambda_2=\kappa_2$, $\lambda_3=\kappa_3$, $\lambda_4=\kappa_4-\tfrac12\kappa_2^2$ and
$\lambda_5=\kappa_5-\kappa_2\kappa_3$ from $\log(1+\sum_{j\ge2}\kappa_j\eps^j)$. Substituting
the values of Theorem \ref{thm:kappavec}, $\lambda_4=\tfrac{146}{29}\zeta(2)^2-\tfrac12\cdot
16\zeta(2)^2=-\tfrac{86}{29}\zeta(2)^2$, and
$\lambda_5=\tfrac{514}{87}\zeta(5)-\tfrac{2200}{87}\zeta(2)\zeta(3)
+4\cdot\tfrac{550}{87}\zeta(2)\zeta(3)=\tfrac{514}{87}\zeta(5)$.
\end{proof}

Since $m=5$ for this operator, $\lambda_5$ is the primitive part of the \emph{first higher}
Frobenius constant. This is the promised occupant of the empty slot (a) of the introduction: no
operator in the literature had its first higher Frobenius constant with a pure
$\QQ\cdot\zeta(5)$ primitive part.

%% FIXED [CRITICAL, same defect]: inherits Theorem \ref{thm:kappavec}'s \VER class.
\begin{observation}[Purity at $j=m$]\label{obs:purityatm}
\PROVED\ from the published values, given the identifications of Theorem \ref{thm:kappavec}
(\VER{$434$ digits}).
For all three operators considered, $\lambda_m\in\QQ\cdot\zeta(m)$ at $j=m$:
\[
  \lambda_2=-\tfrac{7}{5}\zeta(2)\ (m=2),\qquad
  \lambda_3=\tfrac{17}{6}\zeta(3)\ (m=3),\qquad
  \lambda_5=\tfrac{514}{87}\zeta(5)\ (m=5).
\]
Beyond $j=m$ purity breaks: for the Ap\'ery $\zeta(2)$ operator, where $j=5$ is three steps past
$m=2$, $\lambda_5=\zeta(5)-\tfrac15\zeta(2)\zeta(3)$.
\end{observation}

\subsection{Beyond \texorpdfstring{$j=5$}{j=5}: weight inhomogeneity}\label{sec:inhomog}
From $j=6$ on, the Brown--Zudilin $\kappa$-series stops being weight-homogeneous but stays in the
ring of multiple zeta values. This is a positive finding: the Ap\'ery operators, run as controls
with the same code, do \emph{not} do it, their $\kappa_j$ being homogeneous of weight $j$
throughout the published range.

\begin{proposition}\label{prop:lambdahigh}
\VER{434 digits; tolerance $10^{-380}$, coefficient bound $10^{14}$}
\begin{align*}
  \lambda_6 &= \tfrac{38416}{249777}\zeta(3)+\tfrac{272}{3045}\zeta(2)^3
               -\tfrac{392}{7569}\zeta(3)^2,\\
  \lambda_7 &= \tfrac{8835680}{8242641}\zeta(3)-\tfrac{38416}{83259}\zeta(2)^2
               +\tfrac{784}{2523}\zeta(2)^2\zeta(3)-\tfrac{1070}{87}\zeta(7),\\
  \lambda_8 &= \tfrac{1326081904}{272007153}\zeta(3)-\tfrac{8835680}{2747547}\zeta(2)^2
               +\tfrac{1319080}{249777}\zeta(5)+\tfrac{74324}{29435}\zeta(2)^4
               -\tfrac{26920}{7569}\zeta(3)\zeta(5),\\
  \lambda_9 &= \tfrac{4709901251104}{260310845421}\zeta(3)
               -\tfrac{1326081904}{90669051}\zeta(2)^2
               +\tfrac{303388400}{8242641}\zeta(5)-\tfrac{341600}{83259}\zeta(2)^3\\
            &\qquad -\tfrac{1716176}{1975509}\zeta(3)^2
               +\tfrac{48800}{17661}\zeta(2)^3\zeta(3)
               +\tfrac{385264}{1975509}\zeta(3)^3+\tfrac{26920}{2523}\zeta(2)^2\zeta(5)
               +\tfrac{590}{87}\zeta(9).
\end{align*}
\end{proposition}

The mixing pattern is sharp and reproducible:

%% FIXED [CRITICAL]: the "$-3$ recursion" was stated on the MOVING top off-diagonal and is
%% FIXED false on the data of Proposition \ref{prop:lambdahigh} printed immediately above. Exact
%% FIXED rational recheck: c_{j-3}(lam_j)/c_{j-4}(lam_{j-1}) = -3 at j=7, but -3365/294 at j=8
%% FIXED and -3660/4711 (resp. -48158/292755) at j=9. The true law is at FIXED weights:
%% FIXED c_4(lam_j) = -3 c_3(lam_{j-1}) for j=7,8,9, all exactly -3, and
%% FIXED c_5(lam_j) = (-3365/294) c_4(lam_{j-1}) for j=8,9, both exactly -3365/294.
%% FIXED Inherited from work/GAMMA_UNIFICATION.md l.250 and repeated in \S\ref{sec:discussion};
%% FIXED both sites corrected. The weight support, the absent weights and the smoothness claim
%% FIXED were re-checked and are correct as printed.
\begin{observation}[Weight pattern]\label{obs:weights} \VER{$j\le9$}
$\lambda_j$ contains exactly the weights $j$ and $3,4,\dots,j-3$. The weights $j-1$, $j-2$ and
$0,1,2$ never occur. The low-weight coefficients propagate along \emph{fixed} weights with
constant ratios: writing $c_w(\lambda_j)$ for the weight-$w$ coefficient,
\[
  c_4(\lambda_j)=-3\,c_3(\lambda_{j-1})\quad(j=7,8,9),\qquad
  c_5(\lambda_j)=-\tfrac{3365}{294}\,c_4(\lambda_{j-1})\quad(j=8,9),
\]
both exactly. The ratio is \emph{not} constant along the moving top off-diagonal
$c_{j-3}(\lambda_j)/c_{j-4}(\lambda_{j-1})$, which is $-3$ at $j=7$ but $-3365/294$ at $j=8$;
the two statements coincide only at $j=7$. All denominators are
$\{3,5,7,11,29\}$-smooth, with $29$ and $87=3\cdot29$ ubiquitous.
\end{observation}

The mechanism producing this inhomogeneity is \OPEN\ and is, in our judgement, the sharpest
remaining question raised by this computation. It probably requires the polynomial $R(T)$ of
\QUOTED{\cite[Thm.~30]{BV}} for a non-cyclic $\sigma_0$; \cite[Cor.~33]{BV} assumes a single
indicial root, which is not the situation here.

%% ==================================================================
\section{What the \texorpdfstring{$\kappa$}{kappa}-vector is not}\label{sec:exclusions}
%% ==================================================================

Exclusions are reported with the same care as identifications. Every entry below is a statement
about a finite search box.

\subsection{Structural hypotheses}
\EXCL{tolerance $10^{-195}$, coefficient bound $10^{12}$} The following natural shapes for
$\kappa$, for the Brown--Zudilin operator:
\begin{enumerate}[label=(\alph*),leftmargin=2.2em]
\item $\kappa=N\cdot G$ with $N(s)=I(s)/(I_5s^5)$ the rational indicial factor and $G$ weight-graded
--- the natural shape suggested by \cite[Thm.~30]{BV}. The refutation has a clean reason:
$\lambda_2,\dots,\lambda_5$ carry no rational part, whereas $N$ has a nonzero linear coefficient.
\item $\log\kappa$ involving Hurwitz values $\zeta(k,\rho)$ at the other local exponents at
$z=0$ (the hypergeometric shape of \cite[Prop.~26]{BV}).
\item powers of $\log c$.
\end{enumerate}
\EXCL{tolerance $10^{-380}$, coefficient bound $10^{14}$} $\kappa_j$ and $\lambda_j$ in the
\emph{graded} weight-$j$ basis, for $j=6,\dots,16$.

\subsection{The gamma-deformation test}\label{sec:gammadef}
It is natural to ask whether the Brown--Zudilin operator is a gamma-class deformation of a known
Ap\'ery operator, given that their $\kappa_5$'s all lie in the two-dimensional space
$\langle\zeta(5),\zeta(2)\zeta(3)\rangle$. The answer is no, and the reason is instructive.

%% FIXED [CRITICAL, same defect]: the BZ coordinate is Theorem \ref{thm:kappavec}'s, which is
%% FIXED \VER{434 digits}; only the linear algebra is proved.
\begin{observation}\label{obs:cheapplane}
\PROVED\ (the linear algebra), given the identification of $\kappa_5(\mathrm{BZ})$ in
Theorem \ref{thm:kappavec} (\VER{$434$ digits}).
The identity
\[
  174\,\kappa_5(\mathrm{BZ})+987\,\kappa_5(\mathrm{Ap\acute ery}\ \zeta_3)
  -3331\,\kappa_5(\mathrm{Ap\acute ery}\ \zeta_2)=0
\]
holds exactly: in coordinates, $174\cdot\tfrac{514}{87}+987\cdot\tfrac73-3331=0$ and
$174\cdot\bigl(-\tfrac{2200}{87}\bigr)+987\cdot\bigl(-\tfrac{17}{3}\bigr)+3331\cdot 3=0$. It is
also \emph{cheap}: three vectors in a two-dimensional space are always dependent. We record it
because such an identity, found numerically and reported without this remark, would look like a
discovery.
\end{observation}

The content is in the individual identifications, whose $\zeta(2)\zeta(3):\zeta(5)$ ratios are
\[
  -\tfrac{1100}{257}\ (\mathrm{BZ}),\qquad -\tfrac{17}{7}\ (\text{Ap\'ery }\zeta_3),\qquad
  -3\ (\text{Ap\'ery }\zeta_2),
\]
against $+2$ for Brown--Zudilin's top period $2[\zeta(5)+2\zeta(2)\zeta(3)]$.

%% FIXED [MINOR]: \EXCL polarity. Everywhere else the marker is followed by the POSITIVE
%% FIXED proposition that the search excluded; this one was stated as a negation.
\EXCL{tolerance $10^{-193}$, coefficient bound $10^{14}$} A two-term rational relation
between $\kappa_j(\mathrm{BZ})$ and $\kappa_j(\text{Ap\'ery }\zeta_3)$, for $j=5,\dots,13$. (The
hits at $j=2,3,4$ --- $\kappa_2(\mathrm{BZ})=2\kappa_2$, $493\kappa_3(\mathrm{BZ})=1100\kappa_3$,
$58\kappa_4(\mathrm{BZ})=365\kappa_4$ --- are forced, each weight-$j$ space being one-dimensional
for $j\le4$.) The three operators are \emph{not} gamma-deformations of one another in any
two-term sense.

%% FIXED [CRITICAL, same defect]: the BZ row is Theorem \ref{thm:kappavec}'s \VER data.
\begin{observation}[The right statement is the primitive one]\label{obs:primitive}
\PROVED\ (the decomposition $\kappa_5=\lambda_5+\lambda_2\lambda_3$, valid for any log-series);
the BZ row given the identifications of Theorem \ref{thm:kappavec} (\VER{$434$ digits}), the
other two from the published values.
In every case $\kappa_5=\lambda_5+\lambda_2\lambda_3$, with
\[
  (\lambda_2,\lambda_3)=\bigl(-4\zeta_2,\tfrac{550}{87}\zeta_3\bigr),\quad
  \bigl(-2\zeta_2,\tfrac{17}{6}\zeta_3\bigr),\quad
  \bigl(-\tfrac75\zeta_2,2\zeta_3\bigr)
\]
for BZ, Ap\'ery $\zeta(3)$, Ap\'ery $\zeta(2)$ respectively. The space
$\langle\zeta(5),\zeta(2)\zeta(3)\rangle$ is a plane of $(\lambda_5,\lambda_2\lambda_3)$
coordinates, not a plane of new periods.
\end{observation}

\subsection{Reciprocal periods}
\EXCL{tolerance $10^{-193}$, coefficient bound $10^{10}$}
$\kappa_j(\mathrm{BZ})\in\Span_\QQ\{\pi^{-2k}:0\le k\le4\}$ for $j=2,\dots,8$. The lesson of
\S\ref{sec:defect} --- always adjoin $\pi^{-2k}$ to a weight-zero basis --- was applied here and
returned nothing. This is consistent: the reciprocal period lives in the connection-constant
\emph{ratio} $c$, not in $\kappa$, matching \cite[Cor.~31]{BV}, which is a statement about
ratios.

%% ==================================================================
\section{The purity defect \texorpdfstring{$c=-3/\pi^2$}{c = -3/pi^2}}\label{sec:defect}
%% ==================================================================

\subsection{Statement and proof}
The three characteristic rays of \eqref{eq:BZrec} carry, respectively, $Q_n$ (growth
$\lambda_3^n$), the pair $I',I''$ (decay $|\lambda_2|^n$), and the combination $I$ (decay
$\lambda_1^n$). The last is the fastest, and it is the one an irrationality proof would want; but
it is occupied by an impure period. The constant that measures the exile is
\[
  c\;=\;\lim_{n\to\infty}\frac{I''_n}{I'_n},
\]
equivalently the unique scalar for which $\Ihat-cI'$ lands on the minimal ray.

\begin{theorem}\label{thm:defect} \PROVED
\[
  c\;=\;-\frac{1}{2\zeta(2)}\;=\;-\frac{3}{\pi^2}\;=\;
  -0.30396355092701331433163838962918291671307632401674\ldots,
\]
and the minimal-ray class it defines is, up to the factor $4\zeta(2)$, the cellular integral
itself:
\[
  M_n:=\Ihat_n-c\,I'_n=\frac{I_n}{4\zeta(2)} .
\]
\end{theorem}

\begin{proof}
Brown and Zudilin prove, for the whole admissible eight-parameter family and not merely at the
symmetric point, the decomposition $I(\boldsymbol a\cdot n)=2I'(\boldsymbol a\cdot n)+
4\zeta(2)I''(\boldsymbol a\cdot n)$ \QUOTED{\cite[\S2 and \S7]{BZ}}, together with
$|\lambda_1(\boldsymbol a)|<|\lambda_2(\boldsymbol a)|$, checkable at each admissible
$\boldsymbol a$ \QUOTED{\cite[Remark 2]{BZ}}. Let
$V_1=\{u:\limsup_n\log|u_n|/n\le\log|\lambda_1|\}$ be the minimal-ray space; $\dim V_1=1$, and
$I\in V_1$ while $I',\Ihat\notin V_1$. Rewriting the decomposition,
\[
  I_n=4\zeta(2)\Bigl(\Ihat_n+\frac{I'_n}{2\zeta(2)}\Bigr)\in V_1,
\]
so $\Ihat+I'/(2\zeta(2))\in V_1$ already. If also $\Ihat-cI'\in V_1$, subtracting gives
$\bigl(c+1/(2\zeta(2))\bigr)I'\in V_1$; since $I'\notin V_1$ this forces
$c=-1/(2\zeta(2))$. The limit form follows because $\Ihat_n-cI'_n=O(\lambda_1^n\,\mathrm{poly}\,n)$
while $I'_n\asymp\lambda_2^n$ and $|\lambda_1|<|\lambda_2|$.
\end{proof}

\begin{corollary}\label{cor:cone} \PROVED
$c(\boldsymbol a)=-3/\pi^2$ for \emph{every} admissible cone point $\boldsymbol a$. The proof uses
only the family-wide decomposition and $|\lambda_1|<|\lambda_2|$. In particular $c$ carries no
information about $\boldsymbol a$: purity is not deformable within the family.
\end{corollary}

\subsection{The certificate}
Independently of the proof, from the exact ladders ($n\le 360$, exact rational arithmetic) at
$\dps=2600$, with $c^*:=-3/\pi^2$ taken exactly:

\begin{table}[h]
\centering
\caption{Convergence of $\Ihat_n/I'_n$ to $-3/\pi^2$. \VER{441 digits at $n=360$}}
\label{tab:cert}
\begin{tabular}{rrrrrrrr}
\toprule
$n$ & $50$ & $100$ & $150$ & $200$ & $250$ & $300$ & $360$\\
$-\log_{10}\bigl|\Ihat_n/I'_n-c^*\bigr|$ & $61.367$ & $122.713$ & $184.060$ & $245.408$
  & $306.756$ & $368.104$ & $\mathbf{441.722}$\\
\bottomrule
\end{tabular}
\end{table}

%% FIXED [MINOR]: 1.22701 and 1.226963 agree to 4 digits, not exactly.
Agreement grows linearly at slope $1.22701$ digits per $n$. The predicted slope is
$\log_{10}(\lambda_2/\lambda_1)=\log_{10}(16.86411)=1.226963\ldots$, and the two agree to
$4$ digits, so the convergence rate is the predicted $(\lambda_1/\lambda_2)^n$.
This is not a numerical coincidence at $441$
digits; it is the theorem's own error term. \PROVED\ (the slope), \VER{441 digits} (the value).
A third-ray check with $c^*$ exact gives $\log|M_n|/n\to\log|\lambda_1|=-5.29756$ from below with
the standard $n^{\alpha}$ correction, and $I_n=4\zeta(2)M_n$ to full working precision (relative
difference $\sim10^{-2160}$ at $n=360$).

\subsection{Why a 600-digit search missed it}\label{sec:reciprocal}
Earlier exclusions for $c$ --- over a fourteen-term weight-graded basis through weight $8$, over
$\{1,\log2,\text{Catalan},\pi,\pi^2\}$, and others --- were all correct. The basis was simply
missing the inverse period. The constant $c$ has weight $-2$, and an integer-relation search over
a $\QQ$-basis of periods can never see a reciprocal period: $1/\zeta(2)$ lies in the $\QQ$-span
neither of
\[
  \{1,\zeta_2,\zeta_3,\zeta_2^2,\zeta_5,\zeta_2\zeta_3\}
  \qquad\text{nor of the weight-zero ratios}\qquad
  \{1,\tfrac{\zeta_3}{\zeta_5},\tfrac{\zeta_2\zeta_3}{\zeta_5},\tfrac{\zeta_2^2}{\zeta_5},
    \tfrac{\zeta_7}{\zeta_2\zeta_5}\}.
\]
The one test that would have found it, a two-term relation between $c\,\zeta(2)$ and $1$, was
never run.

\begin{remark}[Methodological lesson]\label{rem:lesson}
Adjoin $1/\zeta(2)=6/\pi^2$, and more generally $\pi^{-2k}$, to any weight-zero integer-relation
basis. We state this because it cost a substantial computation, and because it generalises: the
constants attached to \emph{ratios} of connection data are exactly the place where
\cite[Cor.~31]{BV}'s ``periods with $2\pi i$ inverted'' becomes operative.
\end{remark}

\subsection{The Tate reading}
Since $4\zeta(2)=-(2\pi i)^2/6$,
\begin{equation}\label{eq:tate}
  c\;=\;\frac{12}{(2\pi i)^2},
\end{equation}
\PROVED\ (exact). This is precisely the shape predicted by \QUOTED{\cite[Cor.~31]{BV}} for
connection constants of a Picard--Fuchs operator: a period with $2\pi i$ inverted. The defect is
a rational multiple of a negative Tate twist --- weight $-2$, depth $0$, the most degenerate
possible answer. It carries no motivic information beyond the Tate line, and the search for an
exotic new period was searching for something that was never there. The graded picture is:

\begin{table}[h]
\centering
\caption{The three rays of the Brown--Zudilin recurrence. \PROVED}
\label{tab:rays}
\begin{tabular}{llll}
\toprule
ray & eigenvalue & class & period carried \\
\midrule
top $(\lambda_3)$     & $592.079\ldots^{\,n}$   & $Q_n$ & $\QQ(0)$ \\
%% FIXED [MINOR]: the middle eigenvalue is NEGATIVE (roots of 4L^3-2368L^2-188L+1 are
%% FIXED 592.079..., -0.0843843..., 0.00500378...); the minus sign was dropped.
middle $(\lambda_2)$  & $(-0.0843843\ldots)^{\,n}$ & $I',\Ihat$ (span) & $\zeta(5)$ and $\zeta(3)$ separately \\
minimal $(\lambda_1)$ & $0.00500378\ldots^{\,n}$& $I=2I'+4\zeta(2)\Ihat$ & $2\zeta(5)+4\zeta(2)\zeta(3)$ \\
\bottomrule
\end{tabular}
\end{table}

The class exiled from the minimal ray is exiled \emph{by $\zeta(2)$}. This is the same object
responsible for the failure of a weight-$5$ harmonic-monomial decomposition of $P_n$: Brown and
Zudilin's top period $\zeta(5)+2\zeta(2)\zeta(3)$ has depth $2$, and its impurity is what blocks
the elementary decomposition. The rotation $c$ is the $\zeta(2)$-twist and nothing more.

%% ==================================================================
\section{Rate--purity conservation}\label{sec:rpc}
%% ==================================================================

\subsection{The exact gap}
From the exact roots of $4\lambda^3-2368\lambda^2-188\lambda+1$:
\[
\begin{aligned}
  \text{rate of the minimal ray} &= -\log|\lambda_1| = 5.2975613530090625810891422832466528\ldots\\
  \text{rate of the middle ray}  &= -\log|\lambda_2| = 2.4723737227466399605499538672117115\ldots\\
  \text{\emph{purity cost}}\quad \Delta &= \log|\lambda_2/\lambda_1|
    = 2.8251876302624226205391884160349413\ldots
\end{aligned}
\]
\VER{35 digits}. Note $\Delta=\log 16.86411\ldots$ and $\Delta/\log10=1.226963\ldots$, which is
exactly the digits-per-$n$ slope of the $441$-digit certificate of Table \ref{tab:cert}: the
same number, now to $34$ digits. $\Delta$ is the number of nats one must pay to move a linear
form from the middle ray to the minimal ray.

\subsection{The law}
We state the following as a conjecture with explicit falsification conditions, because that is
what it is: a pattern read off three families and then successfully predicted on a fourth.

\begin{conjecture}[Rate--Purity Conservation, (RPC)]\label{conj:rpc}
Let $L$ be an Ap\'ery-like operator, $m$ the multiplicity of the local exponent $0$ in the
indicial polynomial at $z=0$, and $r$ the order of the associated recurrence.
\begin{enumerate}[label=\textup{(RPC-\arabic*)},leftmargin=3.4em]
\item The characteristic rays are indexed by the \emph{depth} $d=0,1,\dots,\lfloor m/2\rfloor$ of
the period they carry; consequently
\[
  \boxed{\;r=1+\lfloor m/2\rfloor\;}
\]
\item Depth and rate are anti-correlated: $|\lambda_{(d)}|$ is strictly decreasing in $d$, and
the price of purifying from depth $d$ to depth $d-1$ is
$\Delta_d=\log\bigl|\lambda_{(d-1)}/\lambda_{(d)}\bigr|>0$.
\item The constant effecting the depth-$1$ to depth-$2$ raise is the middle-ray connection ratio
$c=A_{\Ihat}/A_{I'}$, a Tate class $12/(2\pi i)^2$ of weight $-2$, constant on the whole
admissible cone. Hence purity is not deformable while $\Delta$ is: no motion in the cone converts
rate into purity.
%% FIXED [MAJOR, circular evidence]: the counting half of (RPC-4) is a TAUTOLOGY given (RPC-1)
%% FIXED -- for ANY log-series the depth-<=2 part of [eps^m] exp(sum lambda_j eps^j) has exactly
%% FIXED 1 + #{(a,b): 2<=a<=b, a+b=m} = floor(m/2) monomials, with no computation. So it is
%% FIXED literally (RPC-1) re-expressed. (RPC-4) now retains only its substantive half,
%% FIXED lambda_m in QQ.zeta(m); the counting is demoted to a remark, Table \ref{tab:rpc}'s last
%% FIXED two columns are marked as re-expressions rather than independent checks, and (F2) in
%% FIXED \S\ref{sec:falsify} has been withdrawn as a falsification route.
\item \emph{(Frobenius witness.)} At the nearest conifold, $\lambda_m\in\QQ\cdot\zeta(m)$: the
primitive part of the Frobenius constant at $j=m$ is a rational multiple of $\zeta(m)$, so that
$\kappa_m=\lambda_m+\sum_{a+b=m}\lambda_a\lambda_b+(\text{depth}\ge3)$ has all of its impurity
in the decomposable terms.
\end{enumerate}
\end{conjecture}

\begin{remark}[what (RPC-4) does \emph{not} say]\label{rem:rpc4}
It is tempting to add to (RPC-4) that the number of depth-$\le2$ monomials of $\kappa_m$ equals
$r-1$. That would carry no content beyond (RPC-1). For \emph{any} series
$\log\kappa=\sum_{j\ge2}\lambda_j\eps^j$ whatever, the depth-$\le2$ part of
$[\eps^m]\exp\bigl(\sum_j\lambda_j\eps^j\bigr)$ consists of the single monomial $\lambda_m$
together with the pairs $\lambda_a\lambda_b$, $2\le a\le b$, $a+b=m$, of which there are
$\lfloor m/2\rfloor-1$; so the count is $\lfloor m/2\rfloor$ identically, for every operator and
with no computation. Asserting that it equals $r-1$ is therefore exactly the assertion
$r=1+\lfloor m/2\rfloor$, i.e.\ (RPC-1) again. We record the monomials in
Table \ref{tab:rpc} because they are the concrete shape of the witness, not because they are a
second test.
\end{remark}

Part (RPC-3) is \PROVED\ for the Brown--Zudilin family (Theorem \ref{thm:defect} and Corollary
\ref{cor:cone}). Parts (RPC-1), (RPC-2), (RPC-4) are the conjectural content.

\subsection{Verification}

\begin{table}[t]
\centering
\caption{(RPC) on four Ap\'ery-like families. All multiplicities are exact symbolic
computations. \PROVED\ for each row's $m$ and $r$; the law itself is Conjecture
\ref{conj:rpc}. The last two columns are columns $2$ and $3$ re-expressed
(Remark \ref{rem:rpc4}), not independent checks; the substantive content of (RPC-4) is
$\lambda_m\in\QQ\cdot\zeta(m)$, which is Observation \ref{obs:purityatm}.}
\label{tab:rpc}
\footnotesize
\begin{tabular}{@{}lccccc@{}}
\toprule
family & $m$ & $1+\lfloor m/2\rfloor$ & actual $r$ & depth-$\le2$ monomials of $\kappa_m$ & $r-1$\\
\midrule
Ap\'ery $\zeta(2)$ & $2$ & $2$ & $2$ \checkmark
  & $\lambda_2\ \to\ 1$ & $1$ \checkmark\\
Ap\'ery $\zeta(3)$ & $3$ & $2$ & $2$ \checkmark
  & $\lambda_3\ \to\ 1$ & $1$ \checkmark\\
Brown--Zudilin $\zeta(5)$, order $9$ & $\mathbf 5$ & $3$ & $\mathbf 3$ \checkmark
  & $\lambda_5+\lambda_2\lambda_3\ \to\ 2$ & $2$ \checkmark\\
weight-$7$ cellular, order-$4$ rec. & $\mathbf 7$ & $4$
  & $\mathbf 4$ \checkmark & $\lambda_7+\lambda_2\lambda_5+\lambda_3\lambda_4\ \to\ 3$
  & $3$ \checkmark\\
\bottomrule
\end{tabular}
\smallskip

\footnotesize
Row 1: $D^2-t(11D^2+11D+3)-t^2(D+1)^2$. \quad
Row 2: $D^3-t(34D^3+51D^2+27D+5)+t^2(D+1)^3$. \quad
Row 4: degree $19$.
\end{table}

The fourth row is a genuine prediction. The law was read off $m=2,3,5$; the multiplicity $m$ for
the weight-$7$ operator was then computed for the first time, from the certified order-$4$,
degree-$19$ recurrence of the weight-$7$ family, via $q_j(x)=c_{4-j}(x+j-4)$. Its indicial
polynomial factors exactly as
\begin{multline}\label{eq:z7indicial}
  I(s)=s^{7}\,(s-1)^{3}\,(2s-5)\,\bigl(3690864s^{8}-59053824s^{7}+409061716s^{6}
      -1601726448s^{5}\\
      +3876913912s^{4}-5940014784s^{3}+5627137052s^{2}-3014681200s+699679047\bigr),
\end{multline}
\PROVED\ (exact factorisation over $\QQ$). Hence $m=7$ --- the multiplicity of the exponent $0$
equals the weight, as for $m=2,3,5$ --- and $1+\lfloor7/2\rfloor=4=r$. \checkmark

\begin{observation}[$m$ equals the motivic weight]\label{obs:mweight} \PROVED\ (on four
families.)
In every family examined, the multiplicity of the exponent $0$ at $z=0$ equals the motivic weight
of the family: $m=2,3,5,7$ for weights $2,3,5,7$.
\end{observation}

\subsection{How to falsify}\label{sec:falsify}
A fifth Ap\'ery-like family of weight $w$ falsifies (RPC) if any of the following occurs.
\begin{enumerate}[label=(F\arabic*),leftmargin=2.6em]
\item Its recurrence order satisfies $r\ne 1+\lfloor w/2\rfloor$ (falsifies RPC-1).
\item The $\lambda_w$ of its nearest conifold is not in $\QQ\cdot\zeta(w)$ (falsifies RPC-4).
\item Its middle-ray connection ratio is a non-Tate period (falsifies RPC-3).
\item Its characteristic moduli fail to be strictly ordered by depth (falsifies RPC-2).
\end{enumerate}
By Remark \ref{rem:rpc4} there is no fifth route through the monomial count: that count is a
formal identity and can never fire independently of (F1).
Concrete open predictions: a weight-$4$ family must have $r=3$; a weight-$6$ family $r=4$; a
weight-$9$ family $r=5$ together with $\lambda_9\in\QQ\cdot\zeta(9)$ --- note that the shape
$\kappa_9=\lambda_9+\lambda_2\lambda_7+\lambda_3\lambda_6+\lambda_4\lambda_5+(\text{deeper})$
is itself formal and predicts nothing.
Each of these is a finite computation on any family for which a recurrence is available.

\subsection{Why this is a rate-plus-purity law}
The reason $\zeta(5)$ is out of reach by the Brown--Zudilin family becomes a counting statement.
Weight $5$ admits $\lfloor5/2\rfloor=2$ depths, so the recurrence must have $3$ rays; the pure
forms $I',\Ihat$ are forced onto the depth-$1$ ray, and the depth-$2$ combination
$I=2I'+4\zeta(2)\Ihat$ monopolises the minimal ray. The coefficient $4\zeta(2)$ is
$-(2\pi i)^2/6$, and its inverse $c=12/(2\pi i)^2$ is the cone-constant defect of
\S\ref{sec:defect}. The rate one would need lives on the ray whose period is impure \emph{by
construction of the depth filtration}, and $\Delta=2.8252\ldots$ is the exact toll.

%% ==================================================================
\section{Connection constants}\label{sec:conn}
%% ==================================================================

\subsection{High-precision values}
Exact rational ladders $Q_n,P_n,\Phat_n$ were propagated to $n=520$ from \eqref{eq:BZrec} with
the anchors of \cite{BZ}; the linear forms $I'=Q\zeta(5)-P$, $\Ihat=Q\zeta(3)-\Phat$,
$I=2I'+4\zeta(2)\Ihat$ were formed at $\dps=3200$ (the cancellation in $I$ is
$(\lambda_3/|\lambda_2|)^n=10^{1999}$ at $n=520$, which is why so much working precision is
needed); the Birkhoff--Stokes limit was then taken on each ray at $\dps=400$, $M=200$. All four
exponents came out $\alpha=-5/2$ automatically, as Proposition \ref{prop:conifold} requires.

\begin{table}[h]
\centering
\caption{Connection (Stokes) constants of the Brown--Zudilin family, $u_n\sim A\lambda^n
n^{-5/2}$.}
\label{tab:conn}
\small
\begin{tabular}{@{}llr@{}}
\toprule
constant & value & agreement across $n$\\
\midrule
$A_Q$      & $\phantom{-}0.06667642572715676784165334063934544446962874559\ldots$ & $322$ digits\\
$A_{I'}$   & $-0.75135587647498929950939202170801814994076813034\ldots$          & $288$ digits\\
$A_{\Ihat}$& $\phantom{-}0.22838480022321613498452097882537506442071493158\ldots$ & $288$ digits\\
$A_I$      & $\phantom{-}4.78381719294327891215233374724085460054978566494\ldots$ & $287$ digits\\
\bottomrule
\end{tabular}
\end{table}

As an independent confirmation of Theorem \ref{thm:defect},
$A_{\Ihat}/A_{I'}+3/\pi^2=5.7\times10^{-381}$: the purity defect reconfirmed to $381$ digits by a
method sharing no code with the certificate of Table \ref{tab:cert}. \VER{381 digits}

At the \emph{first} singularity $z_3=1/\lambda_3$ one has $A_P/A_Q=\zeta(5)$ and
$A_{\Phat}/A_Q=\zeta(3)$ \VER{3878 digits}. We flag these as formally trivial --- they say
$\lim P_n/Q_n=\zeta(5)$, which holds because $I'_n\to0$ --- and record them only as a very strong
consistency check on the extrapolation machinery, not as a theorem. Reading the singularities
outward from $z=0$, the connection ratios are
\[
  z_3:\ \zeta(5),\ \zeta(3);\qquad z_2:\ 1/(2\zeta(2));\qquad z_1:\ \text{---},
\]
the weights descending $5$, $3$, and then \emph{inverting} to $-2$. That inversion is the defect.

\subsection{The identifications}
The relation finder was validated first on $\Gamma(1/3)\Gamma(2/3)=2\pi/\sqrt3$, on the
$\Gamma(1/6)$ duplication formula, and on the Ap\'ery $\zeta(3)$ Stokes constant
$S(0)=(1+\sqrt2)^2/(2^{9/4}\pi^{3/2})$, before being applied to unknown targets. Basis hygiene is
discussed in Trap \ref{trap:basis}.

\begin{theorem}[The Stokes determinant]\label{thm:det}
\VER{260 digits; relative difference $1.1\times10^{-260}$}
\[
  A_Q\cdot A_{I'}\cdot A_I \;=\; -\frac{\pi^{5/2}}{12\sqrt{37}},
  \qquad\text{equivalently}\qquad
  \bigl(A_QA_{I'}A_I\bigr)^2=\frac{\pi^{5}}{5328},\quad 5328=2^4\cdot3^2\cdot37 .
\]
Using $c=-3/\pi^2$ this gives, independently verified,
$A_Q\cdot A_{\Ihat}\cdot A_I=\sqrt{\pi}/(4\sqrt{37})$.
\end{theorem}

\begin{theorem}[The individual constants]\label{thm:individual}
Let $\lambda_3=592.0793805346115628\ldots$ be the dominant characteristic root. Then
\begin{equation}\label{eq:AQ}
  64\,\bigl(A_Q\,\pi^{5/2}\bigr)^2=u,\qquad
  37u^3-3219u^2-229u-1=0,\qquad
  u=\frac{-700\lambda_3^2+526604\lambda_3-6199}{762533},
\end{equation}
with $762533=37^2\cdot557$; that is, $A_Q=\sqrt{u}\,/\,(8\pi^{5/2})$ with $u\in\QQ(\lambda_3)$
explicitly. \VER{261 digits} Moreover, writing $t$ for the indicated square,
\[
\begin{aligned}
  t=\bigl(A_{I'}\pi^{-5/2}\bigr)^2&:\quad 1726272\,t^3+4171824\,t^2-8244\,t+1=0,\\
  t=\bigl(A_{I}\pi^{-5/2}\bigr)^2&:\quad 37\,t^3+51504\,t^2-58624\,t+4096=0,\\
  t=\bigl(A_{\Ihat}\pi^{-1/2}\bigr)^2&:\quad 2368\,t^3+51504\,t^2-916\,t+1=0 .
\end{aligned}
\]
\VER{all at tolerance $10^{-230}$, coefficient bound $10^{16}$}
\end{theorem}

Internal consistency, which is a genuine check and not a restatement:
\begin{itemize}[leftmargin=1.6em]
\item The $A_{I'}$ cubic is the $A_{\Ihat}$ cubic under $t\mapsto t/9$ --- forced by
$A_{I'}=A_{\Ihat}/c=-(\pi^2/3)A_{\Ihat}$. \PROVED\ (exact polynomial identity).
%% FIXED [MINOR]: -5/2+5/2+5/2+1/2 = 3, not 5/2. It is the THREE powers of the determinant
%% FIXED triple that sum to 5/2; A_Ihat belongs to the other triple, which sums to 1/2.
\item The three $\pi$-powers $-\tfrac52,+\tfrac52,+\tfrac52$ of $A_Q,A_{I'},A_I$ sum to
$\tfrac52$, matching Theorem \ref{thm:det}; the other triple $A_Q,A_{\Ihat},A_I$ has powers
$-\tfrac52,+\tfrac12,+\tfrac52$ summing to $\tfrac12$, matching $\sqrt\pi/(4\sqrt{37})$.
\PROVED
\item The leading coefficient $2368$ of the $A_{\Ihat}$ cubic is the middle coefficient of the
characteristic polynomial $4\lambda^3-2368\lambda^2-188\lambda+1$; the primes $37$ and $557$ are
those of $41218=2\cdot37\cdot557$; the singular cubic has discriminant $2^4\cdot37^3\cdot557^2$
and $\Disc(37u^3-3219u^2-229u-1)=2^4\cdot37\cdot26357^2$ --- the \emph{same square class} $37$,
i.e.\ the same quadratic resolvent $\QQ(\sqrt{37})$, and indeed $u\in\QQ(\lambda_3)$ explicitly.
\PROVED
\end{itemize}

\subsection{The $\Gamma$-content}
\EXCL{tolerance $10^{-240}$, coefficient bound $10^{6}$} The quantity $\log|A_Q|$ lies in the
$\QQ$-span of
\[
  \{\log\Gamma(k/s)\}\ \cup\
  \{\log\pi,\ \log2,\ \log3,\ \log37,\ \log557,\ \log|z_1|,\ \log|z_2|\},
\]
for $s\in\{3,4,5,6,8,12,24\}$ and combinations thereof.

\EXCL{coefficient bound $10^{20}$} $A_Q\pi^{e}\in\QQ(\lambda_3)$ --- as opposed to its quadratic
extension --- for $e\in\{-3/2,-1/2,0,1/2,3/2\}$.

%% FIXED [MAJOR-4]: the Observation carried no evidence marker, and "no Gamma(r/s) with s>2"
%% FIXED overstated the \EXCL box above, which tested only s in {3,4,5,6,8,12,24} at tolerance
%% FIXED 1e-240, coefficient bound 1e6. The source (GAMMA_UNIFICATION.md l.53-54) keeps the
%% FIXED restriction. Fixed here, in the abstract and in Results item (6).
\begin{observation}[Verdict]\label{obs:gammaverdict}
\EXCL{tolerance $10^{-240}$, coefficient bound $10^{6}$, $s\in\{3,4,5,6,8,12,24\}$}
The only $\Gamma$-value found in the Brown--Zudilin connection constants is
$\Gamma(1/2)=\sqrt{\pi}$. The half-integer power is $\pi^{\rho+1}$ with $\rho=3/2$ the conifold
exponent --- the classical transfer factor $1/\Gamma(-\rho)$, $\Gamma(-3/2)=4\sqrt{\pi}/3$ --- and
the algebraic part is a square root in the $S_3$ cubic field $\QQ(\lambda_3)$ of the singular
locus. Within that box: no $\Gamma(r/s)$ with $s>2$, no cubic-field regulator, no new
transcendental. As Remark \ref{rem:underranged} insists, this is a statement about a box; the
denominators $s\notin\{3,4,5,6,8,12,24\}$ were not searched.
\end{observation}

\begin{remark}[An earlier negative, corrected]\label{rem:underranged}
A previous search reported that $A_Q,A_{I'},A_I$ are ``not algebraic times $\pi^k$'', on the
strength of tests of $A\pi^{e}$ for $e\in\{-\tfrac12,0,\tfrac12,1,\tfrac32,2,\tfrac52,3\}$ with
minimal polynomials of degree $\le12$ and coefficient bound $10^5$. Those exclusions were
correct as stated and simply under-ranged in two ways: the exponents $e=-5/2$ and $e=-1/2$ ---
the right ones for $A_{I'},A_I,A_{\Ihat}$ --- were never tested, and for $A_Q$ the correct
exponent $e=+5/2$ \emph{was} tested, but the minimal polynomial \eqref{eq:AQ} has coefficients up
to $1.3\times10^{7}$, some $130$ times the bound used. We record this because it is the standard
failure mode of exclusion results and is worth stating explicitly: an exclusion is a statement
about a box.
\end{remark}

%% ==================================================================
\section{Application: no pair-worthiness at weight 7}\label{sec:weight7}
%% ==================================================================

We now apply the picture to the question it was built for. Zudilin's theorem that one of
$\zeta(5),\zeta(7),\zeta(9),\zeta(11)$ is irrational \cite{Zu02a,Zu04} remains the narrowest
window containing $\zeta(5)$; a construction producing a small linear form in $1,\zeta(5),\zeta(7)$
alone would narrow it to a pair. The weight-$7$ analogue of the Brown--Zudilin family is the
natural candidate. We show it does not deliver.

\subsection{Setup}
The weight-$7$ family has a certified order-$4$, degree-$19$ recurrence $L$ with palindromic
characteristic polynomial
\[
  7381728\,x^4-46800155520\,x^3+501765579072\,x^2-46800155520\,x+7381728,
\]
whose roots come in reciprocal pairs \PROVED:
\[
  \{6329.2605150\ldots,\ 10.6453896153\ldots,\ 0.0939373791041\ldots,\ 1.5799634059\times10^{-4}\},
\]
with $\log$-rates $\pm8.75293868604055$ and $\pm2.36512689845195$. The ladder consists of
$q_n$, $s_n$, $\Phat_n$ and (unavailable beyond small $n$) $P_n$; the linear forms are
$I''_n=-9q_n\zeta(5)+2s_n\zeta(3)-\Phat_n$ and the primitive form $I'_n$. Write
\[
  \text{sector A}:\ \rho_A=1.5799634\times10^{-4},\qquad
  \text{sector B}:\ \rho_B=0.0939373791041 .
\]
Sector A lies below $e^{-7}=9.1188\times10^{-4}$ and would pass the irrationality criterion;
sector B does not.

\subsection{The sector theorem}
The earlier argument for this family required two facts: that $I''$ rides sector B, and that $I$
rides sector A, so that the sector-B coefficient of $I'=I-\zeta_2 I''$ is nonzero. The second is
a delicate assertion that a cancellation is exact. It is not needed.

\begin{theorem}[Sector exclusion by positivity]\label{thm:sector}
Given that $s,\Phat\in\Sol(L)$ (see \S\ref{sec:residual}), the primitive weight-$7$ form
satisfies
\[
  I'_n \;=\; I_n-\zeta_2 I''_n \;=\; I_n+\zeta_2|I''_n|,\qquad\text{hence}\qquad
  |I'_n|\ \ge\ \zeta_2\,|I''_n| ,
\]
so that $\limsup_n|I'_n|^{1/n}\ \ge\ \rho_B=0.0939374\ldots\ \gg\ \rho_A$. Sector A is
\emph{excluded}, not merely disfavoured.
\end{theorem}

\begin{proof}
$I'$ is a sum of two positive quantities. The two inputs are:
$I''_n<0$ for all $n=0,\dots,140$, verified in exact rational arithmetic \VER{exact, $141$
indices}; and $I_n>0$ for all $n$, which is rigorous, $I$ being an all-positive fourfold series.
Given $s,\Phat\in\Sol(L)$, $|I''_n|^{1/n}$ converges to one of the four discrete values above,
and \S\ref{sec:sectormeas} identifies the limit as $\rho_B$. No conspiracy of coefficients can
make $I'$ decay faster than $I''$.
\end{proof}

The improvement over the cancellation argument is that the verdict is now independent of
everything about $I_n$'s own decay rate, of the denominator of $P_3$, and of any grid hypothesis.
The sole remaining dependency is $s,\Phat\in\Sol(L)$, discussed in \S\ref{sec:residual}.

\subsection{Measuring the sector}\label{sec:sectormeas}
$I''_n$ was evaluated from exact rationals at $\dps=4.9n+100$, so there is no
forward-propagation instability. The statistic is $-\log|I''_{n+1}/I''_n|$; the model
$I''_n\sim C\rho^n n^{\beta}$ gives the raw estimator a $-\beta/n$ bias, killed by Richardson
extrapolation.

\begin{table}[h]
\centering
\caption{Rate of $I''$. Target $-\log\rho_B=2.36512689845$. \VER{5 digits}}
\label{tab:sector}
\begin{tabular}{rrrr}
\toprule
$n$ & raw & Richardson-1 $(n,2n)$ & Richardson-2\\
\midrule
$10$ & $2.6796806457$ & $2.3818974599$ & $2.3655899615$\\
$20$ & $2.5307890528$ & $2.3696668361$ & $2.3651904833$\\
$30$ & $2.4775676766$ & $2.3672005606$ & $\mathbf{2.3651463543}$\\
$40$ & $2.4502279445$ & $2.3663095715$ & ---\\
$50$ & $2.4335830363$ & $2.3658901534$ & ---\\
$60$ & $2.4223841186$ & $2.3656599059$ & ---\\
$69$ & $2.4150358618$ & $\mathbf{2.3655314035}$ & ---\\
\bottomrule
\end{tabular}
\end{table}

This is a \emph{discrete identification, not a continuous fit}: given $s,\Phat\in\Sol(L)$, the only
admissible answers are $\{-8.752939,-2.365127,+2.365127,+8.752939\}$. The Richardson-2 value is
within $2\times10^{-5}$ of one of them and $4.7$ away from the nearest alternative. Two disjoint
windows agree. \textbf{Sector B.}

\subsection{The honest denominator exponent}\label{sec:kappa7}
Whether the true denominator budget is below $7$ is the one remaining lever, and it is not
available. Measured directly on the $361$ exact weight-$5$ ladder terms --- the only family where
the true $P$ is known:
\[
  \den(P_n)\mid 12\,d_n^{5}\ \text{ for } 361/361;\qquad
  \den(P_n)\mid d_n^{5}\ \text{ for } 1/361;\qquad
  \den(P_n)\mid 6\,d_n^{5}\ \text{ for } 1/361,
\]
\VER{exact, $n\le360$}, so the constant $12$ is sharp and the exponent is exactly the weight.
The approach to the exponent is
\[
  \frac{\log\den(P_n)/n}{\log(d_n^5)/n}=0.9953,\ 0.9928,\ 0.9947,\ 0.9936
  \quad\text{at } n=80,200,320,355,
\]
converging to $1$. The weight-$7$ descent ladder behaves identically:
$\log\den(\Phat_n)/n$ divided by $5\log d_n/n$ is $0.9856$ at $n=140$, i.e.\ the descent exponent
$5$ is tight.

\begin{observation}\label{obs:kappa7} \VER{as tabulated}
$\kappa=7$ is the honest denominator budget for the weight-$7$ primitive form, not a placeholder.
There is no $d_n^2d_{2n}$-type structure lowering it: the analogue that can be measured is tight
at its nominal value. (The slack members are $\den(s_n)\sim d_n^{1.95}$ and weight-$5$
$\den(\Phat_n)\sim d_n^{3.95}$, but the binding constraint on a linear form is its constant term,
and that is tight.)
\end{observation}

To pass at sector B one would need $\kappa<2.365$, i.e.\ $\den(P_n)\lesssim d_n^{2.365}$. The
shortfall is $d_n^{4.635}\approx103^{\,n}$.

\subsection{The worthiness numbers}\label{sec:gamma}
With $C_1=\log\lambda_{\max}$, $C_0=\log\rho(I')$, $C_2=\kappa$, the two standard measures are
\[
  \gamma^{\mathrm{ratio}}=\frac{-C_0}{C_2},\qquad
  \gamma^{\mathrm{worth}}=\frac{C_1-C_0}{C_1+C_2},
\]
which exceed $1$ simultaneously, both conditions being equivalent to $-C_0>C_2$. \PROVED. At
weight $5$ the second reproduces Brown--Zudilin's printed symmetric-point worthiness
$\gamma=0.777959763\ldots$ \QUOTED{\cite[\S2]{BZ}} \VER{8 digits}. At weight $7$:

\begin{table}[h]
\centering
\caption{Weight-$7$ worthiness. $C_1=8.75293868604$, $C_2=\kappa=7$. \VER{12 digits}}
\label{tab:gamma}
\begin{tabular}{lrrr}
\toprule
 & $-C_0$ & $\gamma^{\mathrm{ratio}}$ & $\gamma^{\mathrm{worth}}$\\
\midrule
sector B (actual) & $2.36512689845$ & $0.337875271$ & $0.705777240$\\
sector A (hypothetical) & $8.75293868604$ & $1.250419812$ & $1.111276932$\\
\bottomrule
\end{tabular}
\end{table}

The exponent gap is $\kappa-(-C_0)=4.634873$ nats, a factor $e^{4.634873}=103.01$ per step. The
margin is not close.

\subsection{Rhin--Viola leverage, quantified}\label{sec:rv}
The dimensionless invariant that a Rhin--Viola group \cite{RV96,RV01} moves is $(-C_0)/\kappa$.
Measured at weight $5$ over the entire known orbit:
\[
\begin{aligned}
  \text{symmetric point } \boldsymbol a=(1,\dots,1):\quad
    &\frac{-C_0}{\kappa}=\frac{2.47237372}{5}=0.494475,\\[2pt]
  \text{record point } \boldsymbol a=(8,16,10,15,12,16,18,13):\quad
    &\frac{-C_0}{\kappa}=\frac{31.55296935}{49.60574813}=0.636075,
\end{aligned}
\]
\VER{6 digits}, a maximal leverage factor of $1.286364$ ($+28.6\%$). Applying the same relative
leverage to weight $7$ from sector B gives
\[
  \gamma_{5,7}\ \approx\ 0.337875\times1.286364\ =\ 0.434631,
\]
still less than half of $1$. Reaching $\gamma_{5,7}=1$ requires a factor $2.9597$ ($+196\%$).
%% FIXED [MINOR rider]: the factor 6.84 is normalisation-dependent --- on the paper's other
%% FIXED measure gamma^worth the same comparison reads 3.68x. Stated, since the conclusion does
%% FIXED not depend on which measure is used.
In additive terms, the group would buy $0.6773$ nats where $4.6349$ are needed --- a factor
$6.84$ on the $\gamma^{\mathrm{ratio}}$ normalisation, or $3.68$ on $\gamma^{\mathrm{worth}}$;
either way several times the total leverage the Rhin--Viola group is empirically worth
at weight $5$. \VER{$12$ digits}

\begin{corollary}\label{cor:norv}
Asymmetric-orbit optimisation cannot close a sector-B gap. It is worth running only if a
different sector is in play.
\end{corollary}

\subsection{Residual dependency, stated plainly}\label{sec:residual}
Everything in this section is unconditional except the hypothesis $s,\Phat\in\Sol(L)$, which rests
on two things.

\emph{(i) $L$ is the true operator for $q_n$}: empirically certified on $74$ exact terms
($70/70$ relations exactly zero) but not proved --- no creative-telescoping certificate exists.
\VER{exact, 70 relations} \OPEN\ as a proof.

\emph{(ii) An arithmetic certificate.} Propagating $s,\Phat,P,q$ exactly in $\QQ$ gives the
denominator ledger
\[
  \den(s_n)\mid d_n^{5}\ \text{ for } 141/141,\qquad
  \den(\Phat_n)\mid d_n^{5}\ \text{ for } 141/141,\qquad n=0,\dots,140 .
\]
The control experiment propagates three generic rational initial conditions through the same $L$:

\begin{table}[h]
\centering
\caption{Denominator ledger versus generic solutions ($\log_{10}$ of denominators).
\VER{exact}}
\label{tab:ledger}
\begin{tabular}{rrrrrrr}
\toprule
$n$ & $\den(s_n)$ & $\den(\Phat_n)$ & generic-1 & generic-2 & generic-3 & $d_n^5$\\
\midrule
$20$ & $14.5$ & $37.6$ & $63.3$ & $64.5$ & $64.1$ & $41.8$\\
$40$ & $29.0$ & $74.2$ & $116.3$ & $117.6$ & $117.1$ & $78.6$\\
$60$ & $47.2$ & $122.3$ & $180.1$ & $181.4$ & $180.9$ & $124.9$\\
\bottomrule
\end{tabular}
\end{table}

The divisibility $\den\mid d_n^5$ holds $61/61$ for the true ladder and $4/61$ for every generic
solution, the four being trivial small-$n$ cases. A degree-$19$ operator divides by $c_4(n)$ at
every step, so a generic solution drags in the prime factors of $c_4(n)$ for all $n$ and its
denominator overshoots $d_n^5$ by some $55$ digits at $n=60$. That the true $(s,\Phat)$ remain
$d_n^5$-smooth for $141$ consecutive $n$ is evidence for $s,\Phat\in\Sol(L)$ of a strength we
regard as decisive given (i), though it is evidence and not proof.

The same test separates $P$: the $L$-propagated $P$ satisfies $\den(P_n)\nmid d_n^{7}$ for $101$
of the $138$ indices $n=3,\dots,140$, the excess taking only the values $\{1,107,321=3\cdot107\}$,
with $\den(P_n)\mid 321\,d_n^{7}$ for $141/141$. So the $L$-propagated $P$ is arithmetically
clean but wrong: it obeys a bounded-constant $d_n^7$ law, yet the persistent prime $107$ ---
absent from $q,s,\Phat$, which are strictly $d_n^5$-clean --- marks it as a different solution
branch. The exclusion of $P$ from $\Sol(L)$ does not rest on the $107$: the propagated value
$I'_3=43.71$ contradicts $I'_3=I_3+\zeta_2|I''_3|=3.5675\times10^{-5}$ to four significant
figures, the latter following from a rigorous majorant $0<I_3\le3.0903\times10^{-9}$.
\PROVED\ (given the ladder).

\begin{remark}
If $L$ itself were wrong, the entire characteristic-polynomial framework would collapse, sector A
included; so no route to a sector-A pass runs through doubting $L$. We found no evidence
whatsoever pointing toward sector A, and one new structural argument (Theorem \ref{thm:sector})
ruling it out given the ladder.
\end{remark}

%% ==================================================================
\section{Discussion: what the middle-root obstruction means}\label{sec:discussion}
%% ==================================================================

We state what we take the results to establish, and, separately, what they do not.

\subsection{What is established}
\begin{enumerate}[label=(\arabic*),leftmargin=2.4em]
\item \emph{The Brown--Zudilin operator has a well-defined and computable $\kappa$-vector, and
its first higher Frobenius constant has a pure $\zeta(5)$ primitive part.} The dictionary between
irrationality constructions and the gamma-conjecture circle is therefore usable in both
directions at weight $5$, and not only at weights $2$ and $3$.

\item \emph{The purity defect of the family is Tate.} There is no unnamed constant in the middle
of this construction. Theorem \ref{thm:defect} identifies the obstruction's magnitude
($c=12/(2\pi i)^2$) and Corollary \ref{cor:cone} shows it is constant on the entire admissible
cone. This closes a line of search rather than opening one, and we regard closing it as the more
useful outcome: the eight-parameter family cannot be tuned to remove the impurity, because the
impurity does not depend on the parameters.

\item \emph{The obstruction is a counting phenomenon, conjecturally.} If (RPC) holds, then the
number of rays is $1+\lfloor m/2\rfloor$ and the assignment of periods to rays is by depth. In
that case a cellular family of weight $w\ge4$ \emph{must} place its purest forms off the minimal
ray, and the toll $\Delta$ is not a defect of the construction but a structural feature of the
depth filtration. The four-family verification (Table \ref{tab:rpc}) is evidence; the
falsification conditions (F1)--(F4) are the way to test it, and each is a finite computation.

\item \emph{At weight $7$ the same thing happens, and the margin is large.} Theorem
\ref{thm:sector} excludes the fast sector by a positivity argument that needs no cancellation
hypothesis. The measured gap is $4.635$ nats against a Rhin--Viola leverage empirically worth
$0.677$ nats at the same normalisation. The weight-$7$ sector-B verdict is the \emph{generic}
outcome of this construction, not bad luck: weight $5$ shows the identical dichotomy, and there
too the smallest saddle would have proved irrationality while the primitive form provably does not
ride it.
\end{enumerate}

\subsection{What is not established}
The results say nothing about whether $\zeta(5)$ or $\zeta(7)$ is irrational, and nothing about
whether \emph{some other} construction rides the minimal ray with a pure form. (RPC) as stated
constrains cellular and Ap\'ery-like families with a single recurrence and a depth-graded period
assignment; it does not preclude a construction outside that class. Nor do we claim the middle-root
phenomenon is unavoidable in principle: what we claim is that it is stable across the two weights
where it can currently be measured, and that within the eight-parameter cone it cannot be moved.

\subsection{Where the leverage would have to come from}
The measured numbers say where it would not. Orbit optimisation on $\mathcal M_{0,10}$ operates
on $(-C_0)/\kappa$, and at weight $5$ the whole group is worth a factor $1.286$ on that invariant;
weight $7$ needs $2.960$. The productive redirection is therefore not orbit search but a
construction whose \emph{primitive} form rides the small saddle --- that is, an attack on the
phenomenon $C_0=\log|\lambda_{\text{middle}}|$ rather than $\log|\lambda_{\min}|$ itself. Three
concrete handles emerge from the computation.
\begin{enumerate}[label=(\roman*),leftmargin=2.2em]
\item \emph{The inhomogeneity of $\lambda_j(\mathrm{BZ})$ for $j\ge6$}
(Observation \ref{obs:weights}). Its mechanism is \OPEN, and every natural structural hypothesis
we could formulate is excluded (\S\ref{sec:exclusions}). Whatever explains the weight pattern
%% FIXED [CRITICAL knock-on]: the "$-3$ recursion between successive off-diagonal coefficients"
%% FIXED is the false form of Observation \ref{obs:weights}; restated at fixed weight.
$\{j\}\cup\{3,\dots,j-3\}$ and the fixed-weight propagation constants $-3$ and $-3365/294$ is
a genuine feature of a non-MUM Frobenius point and is not present in the rank-$2$ literature.
\item \emph{The second condition of \cite[Def.~21]{BV}} for the Brown--Zudilin operator ---
analyticity at $c$ of the $\sigma_c$-invariant solutions of $L^\vee$ --- is \OPEN, and its
verification would put \S\ref{sec:kappavec} on a fully rigorous footing rather than the empirical
licence of Observation \ref{obs:sufficient}.
\item \emph{A primitive operator} $\Ltilde$ annihilating $I'$ at weight $7$ remains unhuntable by
data fitting: there is no $P$ or $I'$ data to fit beyond $n=3$, and the residue and
high-precision routes are both walled. The one $\Ltilde$-adjacent fact now available is that the
$L$-propagated $P$ satisfies $\den(P_n)/d_n^7\in\{107,321\}$ constantly, so any $\Ltilde$ must
differ from $L$ in a way that changes the constant-term arithmetic without changing $q,s,\Phat$.
\end{enumerate}

\subsection{A remark on evidence}
Three of the findings above began as reported negatives that turned out to be under-ranged or
mis-based: the ``not algebraic times $\pi^k$'' verdict on the connection constants (Remark
\ref{rem:underranged}), the exclusion of $c$ from every period basis tried (\S\ref{sec:reciprocal}),
and a spurious integer relation produced by a degenerate basis (Trap \ref{trap:basis}). In each
case the earlier statement was true of the box that was searched. We have therefore attached to
every exclusion in this paper the tolerance and the coefficient bound, and to every identification
the digit count, and we recommend the practice.

%% ==================================================================
\appendix

\section{Two traps}\label{app:traps}
%% ==================================================================

Both cost real time, and both generalise.

\begin{trap}[Degenerate bases in integer-relation search]\label{trap:basis}
An integer-relation algorithm returns a relation among the vectors it is given. If the basis has
an internal relation, the algorithm will find \emph{that}, with coefficient zero on the target,
and the output looks like a discovery.

Two instances occurred here.
\begin{enumerate}[label=(\alph*),leftmargin=2.2em]
\item \emph{Zeta-value bases.} A weight-$\le8$ basis containing both $\zeta(8)$ and $\zeta(2)^4$
returned the relation $175\zeta(8)=24\zeta(2)^4$ with coefficient $0$ on the target. \PROVED\
(exact: $\zeta(8)=\tfrac{24}{175}\zeta(2)^4$.)
\item \emph{Logarithm bases.} For the Brown--Zudilin singular cubic
$z^3-188z^2-2368z+4$, the product of the roots is $-4$, so
\[
  \log|z_1|+\log|z_2|+\log|z_3|=\log 4
\]
\PROVED\ (exact). Including all three logarithms in a search basis alongside $\log 2$ is therefore
degenerate, and it produced a spurious zero-on-target hit until $\log|z_3|$ was dropped.
\end{enumerate}

\emph{Rule.} Run the relation finder on the basis alone, first. If it returns anything, the basis
is degenerate and any subsequent hit must be re-read.
\end{trap}

\begin{trap}[Resonant Frobenius normalisation]\label{trap:resonance}
The weight-$7$ operator has indicial polynomial \eqref{eq:z7indicial}, which contains the factor
$s^{7}(s-1)^{3}$: the exponents $0$ and $1$ are both present and \emph{resonant}. Consequently
$a_n(s)$ has a triple pole at $s=0$, and the Frobenius deformation is not defined by the naive
normalisation $a_0=1$.

%% FIXED [MINOR]: "not quoted anywhere in this paper" appeared three lines after quoting two of
%% FIXED them. Restated as: not used anywhere, and quoted only here, as the illustration.
Clearing the pole with $a_0(s)=s^3$ produces clean-looking values
$\lambda_2=-\tfrac{292}{55}\zeta(2)$ and $\lambda_3=\tfrac{432}{55}\zeta(3)$, and then
contaminated values from $\lambda_4$ onward: the base point has slid to the exponent-$1$
solution. \emph{The weight-$7$ $\kappa$-values are therefore quarantined: the two displayed in
the previous sentence are quoted here as the illustration of the trap and nowhere else, and no
result of this paper uses any of them.} The (RPC) test on the weight-$7$ family (Table \ref{tab:rpc}, row 4)
rests only on the exact symbolic facts $m=7$ and $r=4$, both of which are independent of the
normalisation. Handling the resonance correctly requires the resonant Frobenius basis and is
\OPEN.
\end{trap}

%% ==================================================================
\section{Computational record}\label{app:record}
%% ==================================================================

All computations were performed in exact rational arithmetic where the statement is exact
(Python \texttt{fractions}, \texttt{sympy}) and in arbitrary-precision floating point otherwise
(\texttt{mpmath} with the \texttt{gmpy2} backend). The following table summarises the run
parameters behind each numerical claim.

\begin{center}
\small
\begin{tabular}{p{0.30\textwidth}p{0.62\textwidth}}
\toprule
claim & parameters \\
\midrule
Instrument validation, \S\ref{sec:validate} &
  $\dps=220$, $K=13$, $M=130$, $n\in\{300,400\}$; self-agreement $208$--$210$ digits;
  residuals $<10^{-219}$ \\
Errata, Observations \ref{obs:err1}--\ref{obs:err2} &
  as above; discrepancies $4.6\times10^{-219}$ and $5.3\times10^{-220}$ against the corrected
  forms \\
$\kappa$-vector, Theorem \ref{thm:kappavec} &
  $\dps=520$, $K=16$, $M=260$, $n\in\{700,800\}$; self-agreement $434$ digits; relation
  tolerance $10^{-380}$, coefficient bound $10^{14}$; residuals $10^{-440}$--$10^{-451}$
  in-session. The archived \texttt{kappas\_hi.json} is written at $420$ digits, so a recheck
  from the archive alone floors at $\approx10^{-419}$ \\
Structural exclusions, \S\ref{sec:exclusions} &
  tolerance $10^{-195}$, bound $10^{12}$ (shapes); tolerance $10^{-380}$, bound $10^{14}$
  (graded bases, $j\le16$) \\
Gamma-deformation exclusion, \S\ref{sec:gammadef} &
  tolerance $10^{-193}$, bound $10^{14}$, $j=5,\dots,13$ \\
Reciprocal-period exclusion &
  tolerance $10^{-193}$, bound $10^{10}$, $j=2,\dots,8$ \\
Defect certificate, Table \ref{tab:cert} &
  exact ladders $n\le360$, $\dps=2600$; $441$ digits at $n=360$; slope $1.22701$ digits per $n$ \\
Connection constants, Table \ref{tab:conn} &
  exact ladders $n\le520$; $I$ formed at $\dps=3200$; Stokes limit at $\dps=400$, $M=200$;
  $287$--$322$ digits \\
Connection identifications, \S\ref{sec:conn} &
  Theorem \ref{thm:det} at relative difference $1.1\times10^{-260}$; \eqref{eq:AQ} at $261$
  digits; cubics at tolerance $10^{-230}$, bound $10^{16}$ \\
$\Gamma$ exclusions, Observation \ref{obs:gammaverdict} &
  $s\in\{3,4,5,6,8,12,24\}$, tolerance $10^{-240}$, bound $10^{6}$; and bound $10^{20}$ for
  $A_Q\pi^{e}\in\QQ(\lambda_3)$ \\
Weight-$7$ ladder, \S\ref{sec:weight7} &
  exact $\QQ$ propagation $n=0,\dots,140$; $I''$ evaluated at $\dps=4.9n+100$; weight-$5$
  denominators on $361$ exact terms \\
\bottomrule
\end{tabular}
\end{center}

%% ==================================================================
\section{Citation discipline}\label{app:cites}
%% ==================================================================

Every reference below was checked against a copy of the source: the bibliographic data were read
from the published text or preprint, and where a value or a statement is attributed to a specific
location in a source, that location was re-read. Three references were added late, and the
following remarks record how they are used:
\begin{itemize}[leftmargin=1.6em]
\item \cite{BeVl25} --- cited only in \S\ref{sec:padic}, for orientation; nothing in this paper
depends on it. The statement attributed to it there is its Theorem 1.4 together with its
equation (2), $\log\Gamma_p(x)=\Gamma_p'(0)x-\sum_{m\ge2}\zeta_p(m)x^m/m$.
\item \cite{BT33} --- for the classical theory of formal solutions at an irregular
singular point of a difference equation. The construction we use is derived from scratch in
Proposition \ref{prop:birkhoff}, so the citation is attributional only.
\item \cite{FBA99} --- for the integer-relation algorithm. All identifications reported
here were re-verified against their closed forms after detection, at the stated precisions, so no
claim depends on the algorithm's provenance.
\end{itemize}

%% ==================================================================
\begin{thebibliography}{99}

\bibitem{Ap79}
R.~Ap\'ery, \emph{Irrationalit\'e de $\zeta(2)$ et $\zeta(3)$},
Ast\'erisque \textbf{61} (1979), 11--13.

\bibitem{Be79}
F.~Beukers, \emph{A note on the irrationality of $\zeta(2)$ and $\zeta(3)$},
Bull. London Math. Soc. \textbf{11} (1979), 268--272.

\bibitem{BV}
S.~Bloch and M.~Vlasenko, \emph{Gamma functions, monodromy and Frobenius constants},
Comm. Number Theory Phys. \textbf{15} (2021), no.~1, 91--147;
\texttt{arXiv:1908.07501}. DOI \texttt{10.4310/CNTP.2021.v15.n1.a3}.

\bibitem{Br16}
F.~Brown, \emph{Irrationality proofs for zeta values, moduli spaces and dinner parties},
Mosc. J. Comb. Number Theory \textbf{6} (2016), no.~2-3, 102--165;
\texttt{arXiv:1412.6508}.

\bibitem{Br26}
F.~Brown, \emph{Mellin transforms, transfinite diameter and rational approximations of
integrals}, preprint (2026); \texttt{arXiv:2604.20741}.

\bibitem{BZ}
F.~Brown and W.~Zudilin, \emph{On cellular rational approximations to $\zeta(5)$},
preprint (2022); \texttt{arXiv:2210.03391}.

\bibitem{Fi04}
S.~Fischler, \emph{Irrationalit\'e de valeurs de z\^eta (d'apr\`es Ap\'ery, Rivoal, \dots)},
S\'eminaire Bourbaki 910, Ast\'erisque \textbf{294} (2004), 27--62;
\texttt{arXiv:math/0303066}.

\bibitem{Fi21}
S.~Fischler, \emph{Linear independence of odd zeta values using Siegel's lemma},
preprint (2021); \texttt{arXiv:2109.10136}.

\bibitem{FSZ}
S.~Fischler, J.~Sprang and W.~Zudilin, \emph{Many odd zeta values are irrational},
Compos. Math. \textbf{155} (2019), no.~5, 938--952; \texttt{arXiv:1803.08905}.

\bibitem{Gol}
V.~Golyshev, \emph{Deresonating a Tate period}, preprint (2009);
\texttt{arXiv:0908.1458}.

\bibitem{GKS}
V.~Golyshev, M.~Kerr and T.~Sasaki, \emph{Ap\'ery extensions},
J. London Math. Soc. \textbf{109} (2024), e12825; \texttt{arXiv:2009.14762}.

\bibitem{GZ1}
V.~Golyshev and D.~Zagier, \emph{Proof of the gamma conjecture for Fano $3$-folds of Picard
rank one}, Izv. Math. \textbf{80} (2016), 24--49. DOI \texttt{10.1070/IM8343}.

\bibitem{GZ2}
V.~Golyshev and D.~Zagier, \emph{Interpolated Ap\'ery numbers, quasiperiods of modular forms,
and motivic gamma functions}, in: Integrability, Quantization, and Geometry II,
Proc. Sympos. Pure Math. \textbf{103.2}, Amer. Math. Soc. (2021), 281--301.

\bibitem{Kerr}
M.~Kerr, \emph{Unipotent extensions and differential equations (after Bloch--Vlasenko)},
preprint (2020); \texttt{arXiv:2008.03618}.

\bibitem{Ko10}
C.~Koutschan, \emph{HolonomicFunctions (user's guide)}, RISC Report 10-01 (2010).

\bibitem{LZ}
L.~Lai and P.~Zhou, \emph{At least two of $\zeta(5),\zeta(7),\dots,\zeta(35)$ are irrational},
preprint (2021); \texttt{arXiv:2103.00904}.

\bibitem{Pa15}
E.~Panzer, \emph{Algorithms for the symbolic integration of hyperlogarithms with applications
to Feynman integrals}, Comput. Phys. Comm. \textbf{188} (2015), 148--166.

\bibitem{RV96}
G.~Rhin and C.~Viola, \emph{On a permutation group related to $\zeta(2)$},
Acta Arith. \textbf{77} (1996), no.~1, 23--56.

\bibitem{RV01}
G.~Rhin and C.~Viola, \emph{The group structure for $\zeta(3)$},
Acta Arith. \textbf{97} (2001), no.~3, 269--293.

\bibitem{Riv}
T.~Rivoal, \emph{La fonction z\^eta de Riemann prend une infinit\'e de valeurs irrationnelles
aux entiers impairs}, C. R. Acad. Sci. Paris \textbf{331} (2000);
\texttt{arXiv:math/0008051}.

\bibitem{BR01}
K.~Ball and T.~Rivoal, \emph{Irrationalit\'e d'une infinit\'e de valeurs de la fonction z\^eta
aux entiers impairs}, Invent. Math. \textbf{146} (2001), no.~1, 193--207.

\bibitem{RV}
B.~Roy and M.~Vlasenko, \emph{Frobenius constants for families of elliptic curves},
Q. J. Math. \textbf{74} (2023), no.~4, 1571--; \texttt{arXiv:2206.15181}.

\bibitem{Zu02}
W.~Zudilin, \emph{A third-order Ap\'ery-like recursion for $\zeta(5)$},
Mat. Zametki \textbf{72} (2002), no.~5, 796--800;
English transl., Math. Notes \textbf{72} (2002), no.~5-6, 733--737.

\bibitem{Zu02a}
W.~Zudilin, \emph{Irrationality of values of the Riemann zeta function},
Izv. Ross. Akad. Nauk Ser. Mat. \textbf{66} (2002), no.~3, 49--102;
English transl., Izv. Math. \textbf{66} (2002), no.~3, 489--542;
\texttt{arXiv:math/0104249}.

\bibitem{Zu04}
W.~Zudilin, \emph{Arithmetic of linear forms involving odd zeta values},
J. Th\'eor. Nombres Bordeaux \textbf{16} (2004), 251--291;
\texttt{arXiv:math/0206176}.

\bibitem{Zu18}
W.~Zudilin, \emph{One of the odd zeta values from $\zeta(5)$ to $\zeta(25)$ is irrational.
By elementary means}, SIGMA \textbf{14} (2018), Art.~028, 8~pp.;
\texttt{arXiv:1801.09895}.

\bibitem{BeVl25}
F.~Beukers and M.~Vlasenko, \emph{Frobenius structure and $p$-adic zeta values},
Adv. Math. \textbf{480} (2025), Part~C, Paper No.~110512;
\texttt{arXiv:2302.09603}. DOI \texttt{10.1016/j.aim.2025.110512}.

\bibitem{BT33}
G.~D.~Birkhoff and W.~J.~Trjitzinsky, \emph{Analytic theory of singular difference equations},
Acta Math. \textbf{60} (1933), 1--89.
\emph{Cited attributionally only --- the construction used here is derived in Proposition
\ref{prop:birkhoff}.}

\bibitem{FBA99}
H.~R.~P.~Ferguson, D.~H.~Bailey and S.~Arno, \emph{Analysis of PSLQ, an integer relation
finding algorithm}, Math. Comp. \textbf{68} (1999), no.~225, 351--369; the algorithm was
introduced in H.~R.~P.~Ferguson and D.~H.~Bailey, \emph{A polynomial time, numerically stable
integer relation algorithm}, RNR Technical Report RNR-91-032, NASA Ames (1992).

\end{thebibliography}

\end{document}
